\chapter{结构化程序设计札记}

{
	\noindent
	\scshape\hfill\scriptsize Edsger W. Dijkstra\hfill
}
\renewcommand{\leftmark}{\normalfont\scriptsize\hfill E. W. DIJKSTRA\hfill}
\medskip

\section[致读者]{致读者}

这些札记可以算是“写给自己的信”：我把它们写下来，是因为如果不这样做，我就会一遍又一遍地重复同样的论证。重读自己写下的文字时，我也并非总是很满意。

其中一个问题是，我觉得它们写得明显过于啰嗦。不过，我（现在）并不打算压缩篇幅。第一，这样做又会耽搁一些时间，而我希望“继续思考”；第二，以往的经历让我担心被误解：许多程序员往往把自己遇到的困难——有时是相当特殊的困难——看成这门学问的核心，于是，对于程序设计究竟是在做什么，人们的看法便大相径庭。

我希望，尽管这些札记有缺点，你仍能喜欢其中至少一部分。如果它们能给你一些启发，或者让你重新认识程序员这门技艺，我的某些目标也就达到了。

在以书的形式出版之前，《结构化程序设计札记》曾私下传阅。读者当时表现出的兴趣，是促使我补充一些材料、让更广泛的读者看到这些札记的主要动力之一；我愿在这里向他们表示感谢。
我尤其要感谢 Bob Floyd、Ralph London 和 Mike Woodger 的鼓励性意见，以及 Peter Naur 提出的批评。最后，我还要感谢 E. L. Dijkstra-Tucker 女士，感谢她在我与英语艰难搏斗时给予的热心帮助。

\section[论我们有限的能力]{论我们有限的能力}
\label{sec:our-inability-to-do-much}

在讲述这些问题时，我面临一个基本困难。我真正关心的是大型程序的构造，其程序文本的篇幅可能与这整本小册子相当。同时，我必须用例子来说明各种技术。出于实际考虑，示范程序必须很小，比我心目中那些“实际规模的程序”小许多倍。我的基本困难恰恰在于：这种规模差异本身，就是程序设计中许多困难的主要来源之一！

如果我能用小型示范程序说明各种技术，然后说：“\dots 而当你面对一个规模大一千倍的程序时，用同样的方法构造它就可以了。”那当然很方便。然而，这种常见的教学手法在这里会适得其反，因为我将阐述的一个中心观点是：任何两种事物，只要在某一方面相差一百倍或更多，就已完全不可比。

历史表明，这个道理很难让人相信。看来，我们受过的训练过于强调忽略规模上的差异，把它们当成“并不本质的程度差异”。我们告诉自己，能做一次的事也能做两次，然后凭着归纳，骗自己相信，只要需要，就能把它做任意多次。但事实并非如此！一千倍这样的差距，已经远远超出了我们的想象能力！

让我举两个例子，把这一点讲得更鲜明些。一个一岁的孩子用手脚爬行，速度可能是每小时一英里。而每小时一千英里则是超音速喷气式飞机的速度。如果把孩子和飞机都看成有移动能力的物体，两者是不可比的：一个能做的事，另一个做不到，反过来也一样。再比如，你可以闭上眼睛，想象自己站在一片开阔地带——草原或海滨——远处一匹没有缰绳的大马正飞奔而来；你能“看见”它跑近，又从身旁跑过。但要在脑海里同样想象一千匹这样的大马列阵奔来，却是不可能的：倘若你真能想象出来，单是恐慌就会让你的心跳停上好几拍！

更麻烦的是，规模问题不仅使我难以讲述，而且正处于这门学问的核心：人们普遍低估规模所带来的特殊困难，看来正是当前软件失败的主要深层原因之一。对此，我只能给出一种回答：尽可能明确地处理规模问题。因此，本节取了这样的标题。

首先是计算的“规模”，也就是计算所涉及的信息量和操作次数。这个规模必须很大；如果它确实很小，索性不用计算机、直接手工完成反倒更容易。自动计算机之所以有存在的理由、之所以有用，恰恰在于它能够完成我们人类无法完成的大规模计算。我们希望计算机做的是我们自己永远做不到的事。今天机器的能力已经强大到这样的地步：即便是小规模计算，仅就其规模而言，也已远远超出我们不借助其他手段时的想象能力。

然而，我们必须把计算组织好，使自己有限的能力足以保证计算产生预期的效果。这种组织工作包括程序构造，于是我们又遇到了另一个规模问题：程序文本的长度。这个问题同样应当得到明确承认。我们应始终记住，读写文本的能力在很大程度上取决于文本的篇幅。在我的国家，电话簿中的条目按城镇或村庄分组，每组内的用户再按姓名的字母顺序排列。我住在一个小村子里；给定一个电话号码，我只需扫过几栏，就能查出它属于谁。但在大城市里做同样的事，就会是一项庞大的数据处理任务！

同样，我也想提醒读者：“清晰”具有显著的数量方面的含义。奇怪的是，许多数学家似乎没有意识到这一点。如果一个定理说，只有满足整整十页的条件，某个结论才成立，那么它就很难成为方便的工具，因为每次引用这个定理，都必须核实所有条件。在欧几里得几何中，勾股定理适用于任意三个点 $A$、$B$ 和 $C$，条件是：可以画出一条经过 $A$ 和 $C$ 的直线，并使它与一条经过 $B$ 和 $C$ 的直线垂直。有多少数学家真正领会到，即便 $A$、$B$、$C$ 中有些点乃至全部点重合，这个定理仍然适用？然而，勾股定理之所以用起来方便，很大程度上似乎正是由于这一点。

总而言之：作为一个头脑迟钝的人，我的脑容量很小。我最好学会接受这一点，尊重并充分正视自己的局限，而不要试图忽略它们；后一种徒劳的努力，将以失败为代价。

\section[论机制的可靠性]{论机制的可靠性}
\label{sec:realiability-of-mechanisms}

我以程序设计为业，谈论的自然是程序，因此本节真正讨论的是程序的可靠性。尽管如此，我还是在标题里用了“机制”一词。这是因为，我把程序看作机制的一种具体形式；而且，我至少想有一次明确表达自己的强烈感受：我关于软件的许多思考，只要作相应调整，同样适用于硬件设计。

今天的计算机是令人惊叹的设备，但最令人惊叹的，却是我们竟然根据如此不确定的理由，认定它们的输出有效。这种不确定性从我们相信硬件运行正常时就已经开始了。

让我们暂且只关注硬件，想一想：我们究竟能在多大程度上确信它的构造正确？几年前，我所在的大学安装了一台计算机。文档中说，除了许多其他部件，它还有一个能将两个 27 位整数相乘的定点乘法电路。于是，一个合理的问题是：“这个乘法器正确吗？它是否按照规格工作？”

对此，一个天真的回答是：“这个乘法器声称能正确完成的不同乘法，数目是有限的，也就是 $2^{54}$；那么，把它们全部试一遍吧。”这个回答看起来合理，实际上却不合理。虽然一次乘法只需几十微秒，但完成这有限的一整组乘法，累计需要的时间却超过一万年！我们必须得出结论：即使只是对乘法器这样的单个部件，穷尽测试也完全不可能。（若以同样的方式测试整台计算机，就意味着必须确认它能正确处理所有可能的程序！）

这一万年的第一个推论是：在整个使用寿命中，乘法器实际被要求完成的乘法，只占全部可能乘法中的一个微不足道的比例，几乎可以说一个也没有！有趣的是，我们仍然要求它在接到指令时，能够正确完成任意一次乘法。这种看似惊人的质量要求，其根本原因是：我们事先不知道，它最终会被要求完成的那极少数乘法到底是哪些。在推理程序时，我们谈论“乘积”，并把因子的具体值抽象掉：我们不知道这些值，也不想知道这些值；知道它们不是我们的职责，不去知道它们才是我们的职责！我们希望用“乘积”这个概念来思考，而不依赖计算中出现的具体实例；这个愿望可以满足，代价却恰恰是这样一项可靠性要求：在庞大的可能集合中，\textit{任何}一次乘法都必须正确完成。以上说明了我们为什么有理由希望得到一个正确的乘法器。

但怎样才能令人信服地确认其正确性呢？只要把乘法器当作黑箱，我们就只能“抽样测试”：向它提供数量上可行的若干对因子，并核对结果。然而，考虑到那一万年的时间，我们显然只能测试全部可能乘法中微不足道的一部分。有些在某种意义上“关键”的乘法，可能整类都没有测试到；相对于我们有理由要求的可靠性，这样的质量控制仍然极不令人满意。因此，实际不会采用这种做法。

由此可以直接得出结论：只要把机制看作黑箱，就无法令人信服地证明它正确；我们唯一的希望，是不再把它看作黑箱。我把这叫作“将机制的结构纳入考虑”。

从现在起，我们所讨论的机制将是程序。（在许多方面，程序比电路更容易处理；电路实际上是一种模拟装置，而且会磨损、老化。）对于程序也是一样：如果不考虑其结构，仅靠测试来确认正确性，做到连最轻微的疑问都排除，也几乎没有希望。换句话说，我们看到：程序正确性能够得到多大程度的确认，并非仅由程序的外部规格和行为决定，而是关键地取决于它的内部结构。

别忘了，我们真正关心的是大型程序。顺便指出，程序的规模本身就要求我们对各个程序部件具有很高的信心。如果单个部件正确的概率为 p，那么由 $N$ 个这样的部件组成的整个程序，其正确的概率大致为
$$
P=p^N.
$$

由于 $N$ 会很大，如果希望 $P$ 与零有显著差别，$p$ 就应当非常、非常接近 $1$！

现在，如果我们认为，程序员不仅有责任写出正确的程序，还有责任令人信服地证明它正确，那么以上讨论就会对程序员的工作产生深刻影响：他所构造的对象，必须具有能发挥作用的结构。

本篇余下的部分，主要探索哪些程序结构可以有效利用。后文将表明，程序正确性并非我唯一关心的事；程序的可适应性或可管理性也是我的关注点。强调程序的可管理性，是我有意作出的选择，因此我想说明这个选择的理由。

过去，普遍可用的设备性能不断增长，减轻了效率要求的紧迫性；但这种增长本身又带来了新的困难。一旦有强大的机器可用，人们就会尝试使用它，所处理问题的规模也会随设备的能力调整：没有人会考虑编写一个需要运行二十年的算法。在过去十到十五年里，处理能力提高了一千倍，人们选择那些现在“技术上应当可行”的问题时，雄心也大大增长。我们想要编写的程序，其规模、复杂程度和精细程度都急剧增加。过去几年中，情况已非常清楚：总体而言，我们的程序设计能力没有跟上这些急剧增长的要求。

可用设备的性能还将继续增长：我们可以预期，厂商会开发更快的机器；即使没有这种发展，今天被认为特别快的机型也会越来越普及。我们想用这些机器做的事，将按相应比例增长。我正是依据这一趋势的外推，形成了自己对程序员任务的看法。

我的结论是：我们迫切需要停止把程序设计主要看成一件让成本$/$性能比最小化的事。
我们应当认识到，即便是现在，程序设计在更大程度上也已经是一项智力挑战：程序设计的艺术，是组织复杂性、驾驭繁多事物，并尽可能有效地避免其衍生的混乱的艺术。

我不愿把效率问题视为程序员首要的关切，并不意味着我忽视它。恰恰相反，效率问题是促使我们修改一个逻辑上正确的程序的主要动力之一。
但我想说的是：只有程序仍然足够易于管理，我们才有余地进行优化——无论所谓优化究竟是什么。

最后，再就计算机的意义说几句题外话。计算机是极其灵活而强大的工具，许多人觉得，它的应用正在改变世界的面貌。我想冒昧地提出这样一个看法：只要仍然主要把计算机当作工具，我们就可能严重低估它的意义。作为工具，它对我们文化的影响或许只是水面上的一圈涟漪；而我预计，它作为一种智力挑战，将产生远为深刻的影响！
\smallskip

\noindent
\textit{本节前半部分的推论：}
\nopagebreak

程序测试可以用来表明错误的存在，却永远不能用来表明错误不存在！

\section[论思维的辅助工具]{论思维的辅助工具}
\label{sec:out-mental-aids}

上一节中，我们指出，程序员有责任让自己的产品“具有能发挥作用的结构”，并且把程序结构与令人信服地证明程序正确性联系起来。

但我们怎样才能说服别人？怎样才能说服自己？哪些典型的思维方式使我们能够理解事物？本节将概略考察这些问题。我先向专业心理学家致以诚挚的歉意，因为这番讨论会带有业余人士的肤浅。不过，我希望——也相信——它足以为我们提供一个尺度，用来衡量某种所提出的结构安排是否有用。

在帮助我们理解程序（或其正确性证明）的思维工具中，有三种我想特别提出来：

\begin{enumerate}[wide, nosep, label=(\arabic*)]
\item 枚举
\item 数学归纳法
\item 抽象。
\end{enumerate}


\subsection{论枚举}
\label{subsec:enumeration}

一组逐一列出的语句按顺序执行，可以引发计算；这组语句中也可以包含用来区分两种或更多情况的条件子句。试图验证这些计算的某项性质，我把这种做法看作借助枚举。下面举一个简单例子，说明我所说的“枚举推理”。

要求证明：以下两条语句依次执行，
\begin{equation*}
	\begin{split}
		&\text{``}dd \coloneq dd/2; \\
		&\text{\textbf{if} } dd \leqslant r \text{ \textbf{do} } r \coloneq r - dd\text{''}
	\end{split}
\end{equation*}

\noindent
在变量“$r$”和“$dd$”上操作时，会使关系
\begin{equation}
	\label{eq:relation-r-dd}
	0 \leqslant r < dd
\end{equation}

\noindent
保持不变。我们只需“跟着”这小段程序走，假定开始时满足式（\ref{eq:relation-r-dd}）。第一条语句把 $dd$ 的值减半，而 $r$ 保持不变；执行之后，关系
\begin{equation}
	\label{eq:relation-r-2dd}
	0 \leqslant r < 2\times dd
\end{equation}

\noindent
成立。现在，分成两种互斥的情况。

（\ref{eq:relation-r-dd}）$dd \leqslant r$。结合式（\ref{eq:relation-r-2dd}），可得关系
\begin{equation}
	\label{eq:relation-dd-r-2dd}
	dd \leqslant r < 2\times dd
\end{equation}

\noindent
在这种情况下，\textbf{do} 后的语句会执行，令 $r$ 减去 $dd$，所以由式（\ref{eq:relation-dd-r-2dd}）可知，最终
\begin{equation*}
	0 \leqslant r < dd,
\end{equation*}

\noindent
也就是说，式（\ref{eq:relation-r-dd}）得到满足。

（\ref{eq:relation-r-2dd}）\textbf{non} $dd \leqslant r$（即 $dd > r$）。在这种情况下，\textbf{do} 后的语句会被跳过，因而 $r$ 也已经是其最终值。此时，“$dd > r$”与第一条语句执行后成立的式（\ref{eq:relation-r-2dd}）结合，立即得到
\begin{equation*}
	0 \leqslant r < dd
\end{equation*}

\noindent
所以，第二种情况下，式（\ref{eq:relation-r-dd}）也得到满足。

这样，我们就完成了关系（\ref{eq:relation-r-dd}）保持不变的证明，也完成了这个包含条件子句的枚举推理例子。

\subsection{论数学归纳法}
\label{subsec:mathematical-induction}

我特意提到数学归纳法，是因为据我所知，它是唯一一种最终能让我们处理循环（例如用重复子句表达的循环）和递归过程的推理方式。我想举一个例子。

考虑数值序列
\begin{equation}
	\label{eq:induction-d}
	d_0,\ d_1,\ d_2,\ d_3,\ \dotso
\end{equation}

\noindent
它由以下关系给定：
\begin{subequations}
\begin{align}
	\label{eq:induction-for-a}
	\text{当 } i &= 0 & d_i &= D \\
	\label{eq:induction-for-b}
	\text{当 } i &> 0 & d_i &= f(d_{i-1})
\end{align}
\end{subequations}

\noindent
其中 $D$ 是给定值，$f$ 是给定的（可计算）函数。要求把变量“$d$”的值置为序列中第一个满足给定（可计算）条件“prop”的值 $d_k$。已知，对于某个有限的 $k$，这样的值存在。更形式化地说，要求建立关系
\begin{equation}
	\label{eq:induction-dk}
	d=d_k
\end{equation}

\noindent
其中 $k$ 由以下表达式（为真）来确定：
\begin{align}
	\label{eq:prod-dk}
	&\text{prop}(d_k)& \\
	\label{eq:non-prop-i}
	\text{且}\quad &\text{\textbf{non} prop}(d_i)& \text{对所有 } i \text{ 满足 } 0 \leqslant i < k
\end{align}

现在考虑下面这段程序：
\begin{equation}
	\label{eq:program-while-non-prop}
	\begin{split}
		&\text{``}d \coloneq D;\\
		&\text{\textbf{while non} prop}(d) \text{ \textbf{do} } d \coloneq f(d)\text{''}
	\end{split}
\end{equation}

\noindent
其中第一行是初始化，第二行是循环，由重复子句 \textbf{while}\dots\textbf{do} 控制；我希望这个子句的含义是不言自明的。（借助上一个例子中使用的条件子句 \textbf{if}\dots\textbf{do}，可以更形式化地定义重复子句的语义，即说明
\medskip

{
	\setlength{\parindent}{8em}
	\hspace{-.4em}``\textbf{while} $B$ \textbf{do} $S$''
}
\medskip

\noindent
在语义上等价于
\medskip

{
	\setlength{\parindent}{8em}
	\hspace{-.4em}``\textbf{if} $B$ \textbf{do}
	
	\quad \textbf{begin} $S$; \textbf{while} $B$ \textbf{do} $S$ \textbf{end}''
}
\medskip

\noindent
这表明，“\textbf{non} $B$”是重复终止的充分必要条件。）

把结构“\textbf{while} $B$ \textbf{do} $S$”中的语句 $S$ 称为“被重复的语句”，我们将证明，在程序（\ref{eq:program-while-non-prop}）中：

\quad 被重复的语句执行第 \textit{n} 次之后，以下关系成立（$n \geqslant 0$）：
\begin{subequations}
	\label{eq:n-execution-prop}
	\begin{align}
		\label{eq:n-execution-prop-a}
		&d=d_n& \\
		\label{eq:n-execution-prop-b}
		\text{且}\quad &\text{\textbf{non} prop }(d_i) &\text{ 对所有 } i \text{ 满足 } 0 \leqslant i < n.
	\end{align}
\end{subequations}

上述断言在 $n = 0$ 时成立（可通过枚举推理确认）；我们必须证明（同样通过枚举推理）：如果它在 $n = N(N \geqslant 0)$ 时成立，那么在 $n = N + 1$ 时也成立。

被重复的语句执行第 $N$ 次之后，关系（\ref{eq:n-execution-prop-a}）和（\ref{eq:n-execution-prop-b}）在 $n = N$ 时得到满足。第 $N + 1$ 次执行发生的充分必要条件是
\begin{center}
	\textbf{non} prop$(d)$
\end{center}

\noindent
为真。根据 $n = N$ 时的式（\ref{eq:n-execution-prop-a}）（即 $d = d_N$），这意味着
\begin{center}
	\textbf{non} prop$(d_N)$
\end{center}

\noindent
从而条件（\ref{eq:n-execution-prop-b}）在 $n = N + 1$ 时也得到满足。此外，由 $d = d_N$ 和式（\ref{eq:induction-for-b}）可得
$$
f(d) = d_{N + 1}
$$

\noindent
因此，被重复的语句
$$
\text{``}d \coloneq f(d)\text{''}
$$

\noindent
第 $N + 1$ 次执行的净效果，是建立关系
$$
d = d_{N + 1}
$$

\noindent
即 $N = N + 1$ 时的关系（\ref{eq:n-execution-prop-a}）；这样，式（\ref{eq:n-execution-prop}）的归纳步骤便得到证明。

现在证明，重复会在被重复的语句执行第 $k$ 次之后终止。当 $n > k$ 时，第 $n$ 次执行不可能发生，因为（根据式 \ref{eq:n-execution-prop-b}）这将意味着
\begin{center}
	\textbf{non} prop$(d_k)$
\end{center}

\noindent
从而违背式（\ref{eq:prod-dk}）。如果重复在被重复的语句执行第 $n$ 次之后终止，那么终止的充分必要条件，即
\begin{center}
	\textbf{non} (\textbf{non} prop$(d)$)
\end{center}

\noindent
借助式（\ref{eq:n-execution-prop-a}），就成为
\begin{equation}
	\label{eq:induction-end}
	\text{prop}(d_n).
\end{equation}

这排除了 $n < k$ 时终止的可能性，因为那会违背式（\ref{eq:non-prop-i}）。所以，重复会在 $n = k$ 时终止。于是，式（\ref{eq:induction-dk}）由式（\ref{eq:n-execution-prop-a}）推出，式（\ref{eq:prod-dk}）由式（\ref{eq:induction-end}）推出，式（\ref{eq:non-prop-i}）由式（\ref{eq:n-execution-prop-b}）推出。我们的证明到此结束。

在结束这个用数学归纳法作为推理方式的例子之前，我想再说几句。因为我有一种不安的感觉：一些读者——特别是有经验、有能力的程序员——读到这里，恐怕已经非常恼火了。对这些读者来说，程序（\ref{eq:program-while-non-prop}）显然正确，所以他们会疑惑，为什么要如此小题大做：“谁都知道，序列中第一个满足某个条件的值是什么意思，他为什么还要像式（\ref{eq:induction-dk}）、（\ref{eq:prod-dk}）和（\ref{eq:non-prop-i}）那样，煞有介事地重述问题？我们还有工作要做，他总不至于期望我们每用一次这样简单的循环，都附上这么冗长、装满数学形式的证明吧？”诸如此类。

老实说，上面这个证明的繁琐和冗长，也让我恼火！但如果确实要证明这段程序正确，目前我还做不到好得多。有时，它引起的愤怒，与我多年前面对平面几何中最初那些简单定理的荒唐证明时一样：那些证明所证明的事情，其“显然程度”和欧几里得的公理本身差不多。

当然，我不敢主张——至少目前不敢！——程序员每在程序里写一个简单循环，就有责任给出这样的证明。如果真如此，他就根本无法写出任何有点规模的程序！这种要求，就像在平面几何中，把每个证明都明确而完整地展开到欧几里得公理一样，不切实际。（参见“\nameref{sec:our-inability-to-do-much}”一节。）

我想说明三点。第一，当程序员觉得式（\ref{eq:program-while-non-prop}）这样的结构显然正确时，是因为他熟悉这种结构。我更愿意把他的行为理解为：他在无意识地引用一个自己\textit{知道}的定理，虽然他可能从未费心把它表述出来；他在一生中的某个时候，曾说服自己相信它为真，虽然他很可能已经忘了自己是怎样做到的，而且当时的方式（很可能）不适合印出来。但我们可以给关于程序（\ref{eq:program-while-non-prop}）的这些断言取个名字，比如“线性搜索定理”。知道这个名字之后，有意识地引用它就更容易，也更自然了。

第二，据我所知，还没有一组与上述例子类似、而且其用途已得到普遍认可的定理。不过，这不足为奇：这样一组定理之所以尚不存在，直接原因是我们所处理的对象类型\textemdash{}也就是程序\textemdash{}还没有定型。程序员所处理的对象，即程序，远没有平面几何所处理的对象那样确定。在此期间，凭直觉胜任工作的程序员，大概会尽可能把自己限定在十分熟悉的程序结构内；而一旦构造出对自己而言不寻常的东西，就变得格外警觉、谨慎。不过，对于一种已经确立的程序设计风格，寻找一批适用于这类程序的定理，可能是一项有益的工作。

第三，最后这个例子所需的证明如此冗长，是一项不应忽视的警告。当然，也许更好的数学家能把证明写得远比我简短、优美。就我个人而言，我倾向于从这种冗长得出结论：程序设计比人们通常认为的更难。让我们诚实地保持谦逊，把证明的长度看作一项迫切的劝告：只要可能，就把自己限定在简单结构内；以智力上的谦逊，像躲避瘟疫一样避开那些“巧妙结构”。

\subsection{论抽象}

在这里，我觉得很难把抽象的作用讲得十分明确，部分原因是它贯穿整个主题。考虑一个算法，以及它可能引发的全部计算：从这些计算出发，把这一次所处理的具体数值抽象掉，剩下的就是算法。“变量”这个概念，代表着对其当前值的抽象。曾有人对我说过（非常遗憾，我记不起是谁，因此无法向应得感谢的人表示感谢）：一个人一旦理解了变量在程序设计中的用法，就理解了程序设计的精髓。再看我们针对重复所运用的数学归纳法，就能为这句话找到印证：一方面，我们借助抽象引入概念，从而得以用这些概念来表述归纳步骤；另一方面，真正要求使用“变量”这个概念的，正是重复。（如果没有重复，就只需使用“量”：它们的值至多定义一次，并不需要像变量的值那样重新定义。）

给一个操作命名，然后根据“它做什么”来使用它，同时完全不管“它如何工作”，这同样涉及抽象。（同理，也应当说，程序设计手册描述的是一台抽象机器：厂商交付的那件具体硬件，只不过是这台抽象机器的机械模型\textemdash{}而且通常并不完美！）在程序中使用一个已命名的操作而不管“它如何工作”，与引用一个定理而不管它如何得到证明，有很强的相似性。即使证明极其复杂，定理本身仍可能用起来十分方便！

这里，我又要提到我们有限的能力。枚举推理在它的能力范围内是好用的，但我们头脑迟钝，这个范围并不大。
只有在一个严格的限制条件下，枚举推理才是合适的思维工具：我们必须十分节制地使用它。
我们应当认识到，抽象是我们减轻枚举推理负担的主要思维手段。

（在这里，英国 Teddington 国家物理实验室的 Mike Woodger 提出以下意见，我怀着谢意把它收入文中：“还有一种平行的类比：表述公理或定理时未加分析的术语，与一个已命名操作预期要处理的、未加分析的操作数之间，也存在这种对应关系。”）

\section[正确性证明的一个例子]{正确性证明的一个例子}

考虑下面这段程序，其中整数常量 $a$ 和 $d$ 满足关系
\medskip

{
	\setlength{\parindent}{8em}
	$a \geqslant 0$\quad 且\quad $d > 0$。
	
	\hspace{-.4em}``\textbf{integer} $r$,\ $dd$;
	
	$r \coloneq a;\ dd \coloneq d$;
	
	\textbf{while} $dd \leqslant r$ \textbf{do} $dd \coloneq 2\times dd;$
	
	\textbf{while} $dd \neq d$ \textbf{do}
	
	\quad \textbf{begin} $dd \coloneq dd/2;$
	
	\quad\quad \textbf{if} $dd \leqslant r$ \textbf{do} $r \coloneq r - dd$
	
	\quad\textbf{end}''.
}
\medskip

为了应用线性搜索定理（参见“\nameref{sec:out-mental-aids}”一节中的“\nameref{subsec:mathematical-induction}”小节），我们考虑以下关系给出的数值序列：
\begin{equation}
	\label{eq:example-initial}
	\begin{split}
		\text{当 } i = 0\qquad & dd_i = d \\
		\text{当 } i > 0\qquad & dd_i = 2\times dd_{i - 1}\\
		\text{由此可得}\qquad & dd_n = d\times 2^n
	\end{split}
\end{equation}


\noindent
其中最后一行可用通常的数学方法推导出来。这些方法也告诉我们，由于 $d > 0$，对于有限的 $r$，
$$
dd_k > r
$$

\noindent
会在某个有限的 $k$ 上成立，从而保证第一个重复会终止，终止时
$$
dd = d\times 2^k
$$

把关系
$$
d_i = 2\times d_{i-1}
$$

\noindent
改写为求 $d_{i-1}$ 的形式，得到
$$
d_{i-1} = d_i/2
$$

\noindent
于是，线性搜索定理告诉我们，第二个重复也会终止。（事实上，第二个被重复语句的执行次数，与第一个完全相同。）

第一个重复终止时，
$$
dd=dd_k
$$

\noindent
因此，
\begin{equation}
	\label{eq:example-linear}
	0 \leqslant r <dd
\end{equation}

\noindent
成立。正如前文所示（“\nameref{sec:out-mental-aids}”一节中的“\nameref{subsec:enumeration}”小节），第二个子句中被重复的语句使这一关系保持不变。重复终止后，根据“\textbf{while} $dd \neq d$ \textbf{do}”，可得
$$
dd = d
$$

\noindent
与式（\ref{eq:example-linear}）结合，得到
\begin{equation}
	\label{eq:example-r-interval}
	0 \leqslant r < d
\end{equation}

此外，我们证明，初始化之后，
\begin{equation}
	\label{eq:example-dd-equiv-0-mod-d}
	dd \equiv 0 \bmod d
\end{equation}

\noindent
成立。例如，这可以从 $dd$ 的可能取值推出（参见式（\ref{eq:example-initial}））：
$$
d\times 2^i\quad \text{对于}\quad 0 \leqslant i \leqslant k.
$$

下一步，我们验证：对 $r$ 完成初始赋值之后，关系
\begin{equation}
	\label{eq:example-a-equiv-r-mod-d}
	a \equiv r \bmod d
\end{equation}

\noindent
始终成立。

\begin{enumerate}[wide, nosep, label=(\arabic*)]
	\item 初始赋值之后，它成立。
	\item 第一个子句中被重复的语句（“$dd \coloneq 2\times dd$”）维持式（\ref{eq:example-a-equiv-r-mod-d}）的不变性，因此整个第一个重复也维持式（\ref{eq:example-a-equiv-r-mod-d}）成立。
	\item 第二个被重复的语句由两条语句组成。第一条（“$dd \coloneq dd/2$”）使式（\ref{eq:example-a-equiv-r-mod-d}）保持不变；第二条也使式（\ref{eq:example-a-equiv-r-mod-d}）保持不变，因为这条语句要么不改变 $r$，要么从 $r$ 中减去 $dd$ 的当前值。根据式（\ref{eq:example-dd-equiv-0-mod-d}），后一操作也维持式（\ref{eq:example-a-equiv-r-mod-d}）成立。因此，整个第二个被重复的语句使式（\ref{eq:example-a-equiv-r-mod-d}）保持不变，从而整个重复也使式（\ref{eq:example-a-equiv-r-mod-d}）保持不变。把式（\ref{eq:example-r-interval}）与式（\ref{eq:example-a-equiv-r-mod-d}）结合，可知最终值满足
\end{enumerate}
\medskip
$$
0 \leqslant r < d\quad \text{且}\quad a \equiv r \bmod d
$$

\noindent
即 $r$ 是 $a$ 除以 $d$ 的最小非负余数。

\begin{enumerate}[wide, nosep, label=\textit{说明 }\arabic*.]
	\item 程序
	
	{
		\setlength{\parindent}{8em}
		\hspace{-.4em}``\textbf{integer} $r,\ dd,\ q;$
		
		$r:= a;\ dd:= d;\ q:= 0;$
		
		\textbf{while} $dd \leqslant r$ \textbf{do} $dd \coloneq 2 \times dd;$
		
		\textbf{while} $dd \neq d$ \textbf{do}
		
		\quad\textbf{begin} $dd \coloneq dd/2;\ q: = 2 \times q;$
		
		\quad\quad\textbf{if} $dd \leqslant r$ \textbf{do begin} $r \coloneq r - dd;\ q: = q + 1$ \textbf{end}
		
		\quad\textbf{end}''
	}
	
\noindent
把对应的商赋给 $q$。观察关系
	$$
	a= q \times dd + r.
	$$
	的不变性，就能完成证明。
	
	（这个例子得自我的同事 N. G. de Bruijn。）

	\item 在“\nameref{subsec:mathematical-induction}”小节中，我们证明了线性搜索定理。在刚才的证明里，又使用了一个关于重复的定理（显然，这个定理只能用数学归纳法证明，不过证明很简单，就留给读者作为练习）：如果进入一个重复之前，某个关系 $P$ 成立，而被重复的语句单次执行不会破坏其成立，那么重复终止后，关系 $P$ 仍然成立。这是一个非常有用的定理，常常让我们不必显式地使用数学归纳法。（这个定理还可以表述得更精确些：对于重复
	$$
		\text{``\textbf{while} } B \text{ \textbf{do} } S\text{''}
	$$
	
	\noindent
	我们必须证明 $S$ 具有如下性质：在执行 $S$ 之前，
	$$
	P \text{ \textbf{and} } B
	$$
	
	\noindent
	为真，能够推出执行之后
	$$ P $$
	
	\noindent
	为真。）
	\item 给读者一道练习（这道题由美国 Pittsburgh 的 CMU 的 James King 提供，谨致谢意）：证明，对于整数 $A$、$B$、$x$、$y$ 和 $z$，在
	$$
	A > 0\quad \text{且}\quad B\geqslant 0
	$$
	
	\noindent
	的条件下，执行程序段
	
	{
		\setlength{\parindent}{8em}
		\hspace{-.4em}``$x \coloneq A;\ y \coloneq B;\ z \coloneq 1;$
		
		\textbf{while} $y \neq 0$ \textbf{do}
		
		\quad\textbf{begin if} odd($y$) \textbf{do begin} $y \coloneq y - 1; z \coloneq z \times x$ \textbf{end};
		
		\quad\quad $y \coloneq y/2;\ x \coloneq x \times x$
		
		\quad\textbf{end}''
	}
	
	\noindent
	之后，最终 $z = A^B$ 成立。
\end{enumerate}

证明必须表明，尽管存在“$y: = y/2$”，所有变量仍保持整数值；证明方法是说明以下关系的不变性：
$$
x > 0\quad \text{\textbf{and}}\quad y \geqslant 0\quad \text{\textbf{and}}\quad A^B = z \times x^y
$$

\section[论证明的有效性与实现的有效性]{论证明的有效性与实现的有效性}

上一节中，我假定了“完善的算术”。根据我的经验，人们常常质疑这类证明的有效性，理由是在实践中，我们从来无法使用完善的算术：允许的整数值通常有一个绝对上限，实数的表示精度也总是有限的，等等。那么，这些证明到底有什么效力？

对此，答案似乎如下：如果证明一个程序正确时，假定的是一个理想化、完善的世界，那么，当这个理想程序交由一个“不完善”的实现执行而出了问题时，就不该觉得奇怪。这是显而易见的！因此，如果想在更现实的世界中证明程序正确，做法就是从一开始便承认：程序中所使用的全部操作（特别是全部算术操作）不一定完善；但我们要以近乎公理化的方式，说明它们为正确执行程序而必须满足的性质，也就是正确性证明所依赖的性质。（上一节的例子中，这项要求只是：在 $[0, 2a]$ 范围内进行精确的整数算术。）

当编写一个对实数进行舍入运算的程序时，必须清楚自己所作的假定，例如
\smallskip

\tabto*{10em}$b > 0\quad \text{蕴含}\quad a + b \geqslant a$

\tabto*{10em}$a\times b =\ b\times a$

\tabto*{10em}$-(a \times  b) =\ (-a)\times b$

\tabto*{10em}$0\times  x =\ 0$

\tabto*{10em}$0 + x =\ x$

\tabto*{10em}$1\times x =\ x$\quad 等等。
\smallskip

这些关系成立，往往是程序逻辑的必要条件。为了兼容性，程序员最好尽可能少提要求；而一个良好的实现，则应尽可能满足更多合理的要求。

在这里，我要坦白自己的一个失误。实现 ALGOL 60 时，我们决定：不仅在两值精确相等时，“$x = y$”返回 \textbf{true}；如果两个值只在所表示的最低有效位上不同，也返回 true。否则，算出 \textit{true} 的可能性实在太小。我们当时想到的是收敛迭代，它可能在舍入精度范围内振荡。尽管我们（出于最好的意图！）宽松地认定实数相等，却很快发现，所选操作的性质太弱，几乎毫无用处。归根结底，确认 $a = b$ \textbf{and} $b = c$ 为真，并不能让程序员推出 $a = c$ 为真。这个决定很快就改了。正是这段经历让我明白，程序员只有凭借工具的（若干）性质，才能使用它；反过来，程序员必须能够说明，自己要求哪些性质。（程序员通常不这样做，因为对于哪些性质可以视为当然，尚无一套约定俗成的传统；如果逐一说明，就会比原本所希望的详尽程度更高。大量配备糟糕浮点运算硬件的机器，加上把自动计算机主要看作数值分析人员工具的误解，给这个行业带来了很大损害！）

\section[论理解程序]{论理解程序}
\label{sec:undertanding-programs}

我见过许多程序设计课程，实质上和通常的驾驶课一样：教的是怎样操纵汽车，而不是怎样用汽车到达目的地。

我想说的是：程序从来不是目的本身。程序的目的在于引发计算，而计算的目的在于产生所期望的效果。虽然程序是程序员制作的最终产品，但程序可能引发的那些计算\textemdash{}它们的“制作”留给机器来完成！\textemdash{}才是这门技艺真正处理的对象。例如，每当程序员说自己的程序正确时，他实际上是在对程序可能引发的计算作出断言。

整个活动的最后阶段，也就是从（静态的）程序文本转变为（动态的）计算，实质上留给机器完成；这一事实增加了困难。从某种意义上说，构造程序因此比建立数学理论更难：程序和理论都是具有结构、不随时间变化的对象。但数学理论本身就有意义，程序却只有通过执行才有意义。

本节余下部分只讨论为顺序机器编写的程序，并探索这样一项责任所带来的若干后果：我们必须借助对程序的理解，对随后发生的计算作出断言。我提出一个（未经证明的）主张：这件事能做得多容易、多可靠，关键地取决于两者之间的关系有多简单，尤其取决于执行顺序控制的性质。笼统地说，我们希望程序文本的结构反映计算的结构。换一种说法：“静态的程序文本在‘文本空间’中展开，而对应的计算在时间中演进；我们能做些什么来缩小两者之间的概念距离？”

计算的目的是产生某种期望的效果。假定它在离散时刻 $t_0$ 开始，在稍后的离散时刻 $t_1$ 完成；我们还假定，可以通过比较“$t_0$ 时的状态”和“$t_1$ 时的状态”来描述它的效果。如果不考虑任何中间状态，就把这一效果视为由一个原子动作产生。

如果考虑若干中间状态，就意味着我们把发生在时间中的事情作了分解。我们把它看成顺序计算，即若干子动作在时间上的相继发生；而且必须确信，这些子动作依次产生的累积效果，确实等于整个计算所期望的净效果。

最简单的情况，是把它分解成固定数量、可以逐一列出的子动作。用流程图表示如下。

\centerline{\includegraphics[scale=0.4]{chap1/fig1}}

这种分解的有效性，必须通过枚举推理确认。在这种情况下，要缩小程序与计算之间的概念距离，可以要求：一段线性排列的程序文本，按照子动作应当发生的次序，写出它们的名称或描述。前面的例子（$0 \leqslant r < dd$ 的不变性）
\medskip

{
	\setlength{\parindent}{8em}
	\hspace{-.4em}``$dd \coloneq dd/2;$
	
	\textbf{if} $dd \leqslant r$ \textbf{do} $r \coloneq r - dd$''
}
\medskip

\noindent
就满足这个条件。对计算的初层分解，是两个动作在时间上的相继执行；在程序文本中，我们能认出这个结构：
\medskip

{
	\setlength{\parindent}{8em}
	\hspace{-.4em}``将 $dd$ 减半；
	
	求 $r$ 对 $dd$ 的余数''。
}
\medskip

我们考虑的是全部满足 $0 \leqslant r < dd$ 的初始状态；对由此考虑的所有计算，以上分为两个子动作的分解都适用。到这里，一切都好。

不过，编写这段程序时假定：“求 $r$ 对 $dd$ 的余数”不是原子动作，而“从 $r$ 中减去 $dd$”是原子动作。考察“求 $r$ 对 $dd$ 的余数”期间所有可能发生的情况时，就有必要作出区分：有些情况下，会执行“从 $r$ 中减去 $dd$”；另一些情况下，$r$ 保持不变。写成
\medskip

{
	\setlength{\parindent}{8em}
	\hspace{-.4em}``\textbf{if} $dd \leqslant r$ \textbf{do} 从 $r$ 中减去 $dd$''
}
\medskip

\noindent
就表达了：在当前这个细节层次上，“求 $r$ 对 $dd$ 的余数”可以采取两种互斥形式之一；同时也给出了在两者之间作选择的依据。如果把“\textbf{if} $dd \leqslant r$ \textbf{do}”看作附加在“从 $r$ 中减去 $dd$”上的条件子句，那么把条件子句放在受它控制的语句前面，自然是合适的。（从这个意义上说，分支子句
\medskip

{
	\setlength{\parindent}{8em}
	\hspace{-.4em}``\textbf{if} 条件 \textbf{then} 语句 1 \textbf{else} 语句 2''
}
\medskip

\noindent
对“语句 1”和“语句 2”施加了“多余的次序”：它们仅仅是两个备选项，只是无法在线性媒介上同时表达。）

分支子句已经由 C. A. R. Hoare 推广；他的“case-of”结构可以在两个以上的可能项之间作选择。它们的流程图表示如下。

\centerline{\includegraphics[height=11em]{chap1/fig2a}\qquad \includegraphics[height=11em]{chap1/fig2b}}

\centerline{\includegraphics[width=.9\textwidth]{chap1/fig3}}

这些流程图具有一个共同性质：顶部只有一个入口，底部只有一个出口。正如虚线框所示，忽略虚线内的内容，就又可以把它们理解为顺序计算中的单个动作。更精确地说，我们处理的是大量可能的计算；对它们的初层分解，都得到相同的子动作时间序列。只有进一步检查\textemdash{}也就是查看虚线框内部\textemdash{}才会发现：在这组可能的计算中，这样一个子动作可以采取某个已列举的不同形式集合中的一种形式。

上述结构足以处理这样一类计算：它们在初层分解中，都分成同一组逐一列出的子动作。但对另一类计算，它们还不够：这类计算在初层分解中，子动作的数目会变化（即在所考虑的这一类计算之间变化）。这时，重复子句的用处就显现出来了。我们提到“\textbf{while} 条件 \textbf{do} 语句”和“\textbf{repeat} 语句 \textbf{until} 条件”，其流程图表示如下。

\centerline{\includegraphics[height=11em]{chap1/fig4a}\qquad \includegraphics[height=11em]{chap1/fig4b}}

这些流程图同样具有顶部单一入口、底部单一出口的性质。它们让我们能够表达：虚线框所代表的动作，进一步检查后，是某一类型的子动作按时间顺序相继执行“足够多次”。

现在，我们已见到三种分解方式，分别可以叫作“串接”“选择”和“重复”。前两种靠枚举推理理解，最后一种靠数学归纳法理解。

如果程序仅用选择子句和重复子句控制执行顺序，那么它可以直接翻译到另一种程序设计语言中：该语言在其他方面都相同，只是执行顺序控制必须通过跳转到带标号的位置来表达。反过来却不成立。换句话说：与可以从任意程序块画箭头指向任意其他块的流程图相比，把自己限定在上述三种分解方式内，所得到流程图的拓扑结构受到限制。相对于那种更大的自由，只使用这些子句，就表现为一种执行顺序上的纪律。

我为什么建议遵守这种执行顺序纪律？这个决定的理由可以从许多方面说明。让我尝试几种说法，希望至少有一种能打动读者。

我们的最终目标之一，是让程序具有良好的结构，使理解程序所需的智力努力（以某种粗略意义衡量）与程序的长度（以另一种同样粗略的意义衡量）成正比。尤其要防止枚举推理的负担急剧增长。这项任务迫使我们采用古老格言“分而治之”的某种做法，也正是我们建议逐步分解计算的原因。

串接式分解可以通过枚举推理来理解。（前提是：在初层分解中，计算被拆分的子动作数量足够少，且对它们净效果的规格说明足够简明。我将在后面再谈这些要求；目前先假定条件已满足。）这样，根据程序文本对计算作出断言就是可行的，因为计算中的进展与程序文本中的进展之间，关系极其简单。特别是：如果进一步检查发现，某个子动作由选择子句或重复子句控制，这也不会加重理解初层分解的负担，因为在那一层，起作用的只有这个子动作的净效果。

由此可得：如果进一步检查发现某个子动作由选择子句控制，那么它具体走了哪条路径，在初层分解中始终无关紧要（唯一重要的是，它走了正确的路径）。同样，如果进一步检查发现某个子动作由重复子句控制，那么被重复的语句执行了多少次，这个数目本身也无关紧要（唯一重要的是，它重复了正确的次数）。

选择子句本身，也可以通过枚举推理理解；重复子句，则可以通过数学归纳法理解。对这三种分解方式\textemdash{}在我看来，这是一大帮助\textemdash{}我们都知道适用的推理模式。

所提出的执行顺序纪律，还有一个好处。理解程序时，我们要建立各种关系。在枚举推理的例子中，我们确认，程序段
\medskip

{
	\setlength{\parindent}{8em}
	\hspace{-.4em}``$dd: = dd/2;$
	
	\textbf{if} $dd \leqslant r$ \textbf{do} $r \coloneq r - dd$''
}
\medskip

\noindent
使关系
$$
0 \leqslant r < dd
$$

\noindent
保持不变。然而，即使我们能保证执行所引程序段之前这些关系成立，也不能断言它们始终成立；具体说，在所引的两条语句执行之间，它们未必成立。换言之，这些关系是否成立，取决于计算进行到了哪里；这似乎是顺序过程的典型特征。

类似地，我们还会赋予变量某种含义：一个变量可能用于计数某类事件发生的次数，比如当前页已经打印的行数。转到下一页之后，立即将它重置为零；打印一行之后，立即加 1。同样，在重置或递增这个计数之前，“当前页已经打印的行数”这一解释是不成立的。因此，赋予变量这样的含义，也只能相对于计算的进展来做。这一观察引出一个问题：“怎样刻画计算的进展？”

简而言之，我们要找一个坐标系，以便标识计算进展中的离散位置；而且希望这个坐标系独立于那些在程序控制下被操作的变量。如果需要用这些变量的值来描述计算进展，我们就陷入了循环论证，因为我们恰恰是想相对于这种进展来解释这些变量的含义。

（还有一个更强的理由说明不应依赖变量的值：设某个程序含有一个不终止的循环，它在有限个不同状态之间反复运行。永远循环的原因在于，计算进展中的不同位置可以处于\textit{同一个}状态。如此一来，状态显然不能区分这两个\textit{不同的}进展位置！）

我们可以换一种方式说明这个问题。给定一个正在运行的程序，假定计算尚未完成，就在某个离散的进展位置停了下来。我们怎样标识这个中断位置？例如，若想把计算重新做到完全相同的位置，该怎么办？又或者，如果停止是由某种运行时错误引起的，我们怎样才能在不完整转储存储器的情况下，标识当时的进展位置？

为简单起见，假定程序文本在（线性的）文本空间中展开，并假定有一种标识机制，用来标识与计算进展的离散位置相对应的程序位置；我们把这种机制叫作“文本索引”。（如果计算进展的离散位置位于相邻两次语句执行之间，那么文本索引所标识的，例如可以是分号。）文本索引是一种广义的指令计数器，其值指向文本中的一个位置。

如果只采用串接和选择进行分解，那么单个文本索引就足以标识计算进展。加入重复子句之后，文本索引就不足以描述计算进展了。不过，每次进入一个重复子句，系统可以引入一个所谓的“动态索引”，严格计数当前相应重复的序号；重复终止时，系统再移除对应的动态索引。由于重复子句可以相互嵌套，合适的机制是栈（即后进先出的存储器）。最初栈为空；进入重复子句时，把一个新的动态索引（置为零或一）放到栈顶；每当决定继续重复，就将栈顶元素加 1；每当决定结束重复，就移除栈顶元素。（这种安排十分清楚地反映了：重复终止之后，重复的次数，甚至它曾经是一个重复这一事实，都不再重要。）

一旦程序设计语言允许使用过程，单个文本索引就不够了。如果文本索引指向某个过程体内部，那么只有同时说明指的是该过程的哪一次调用，才能刻画计算的动态进展；而这可以通过给出一个指向调用位置的文本索引来做到。引入过程后，文本索引必须推广为文本索引栈：调用过程时增加一个元素，过程返回时减少一个元素。

关键在于，这些索引的值不受程序员控制。不论他是否愿意，它们都由程序的书写形式或当前计算的动态演进所定义。它们提供了一套独立坐标，用来描述计算进展；在这个“独立于变量”的参照系中，可以给变量赋予含义。

当然，即便自由使用跳转，也存在一个独立于程序员的坐标系，可以用来唯一描述顺序计算的进展：一种规范化的时钟，计数从程序开始起经过了多少个“计算进展的离散位置”。它确实能唯一标识位置，却完全无助于理解，因为文本索引已不再是这个坐标系的组成部分。

这里的教训是：如果承认自己有责任借助引发计算的程序文本，在智力上掌握这些计算，那么我们就应当谦逊地把自己限定在最系统的执行顺序机制内，确保“计算中的进展”以最直接的方式映射到“文本中的进展”。

\section[论程序的比较]{论程序的比较}
\label{sec:comparing-programs}

程序员在日常工作中经常遇到这样的情况：对于一个给定问题，即使已指定用某个算法求解，又指定了机器，程序仍远非唯一确定。设计过程中，他必须在不同方案间选择；得到一个正确程序之后，也经常需要修改它。例如，人们可能觉得，另一个方案的计算对现有设备资源的要求更合适。

这些情况引出了程序等价性的问题：给定两个程序，它们是否引发净效果相同的计算？如果对程序的给定方式、执行这些计算的机器，以及计算的“净效果”都作适当形式化，这大概可以成为一个表述明确、能吸引某些数学头脑的问题。但我不打算以这种一般形式处理它。恰恰相反，我并不是从两个任意给定的程序出发——例如由两个作者独立构思的程序；我关心的是那些可以视为同一个人思考结果的不同程序方案。于是，问题变成：怎样构思这两个方案，并安排它们的结构，才能使比较工作更容易？

我做过许多试验，从中得到的基本经验可以概括如下：两个程序若引发净效果相同的计算，它们便\textit{在这个意义上}等价，而\textit{不能事先认定}它们在其他任何意义上也等价。当我们想通过比较程序，来比较其对应的计算时，基本经验是：如果在所比较的层次上，两个程序中的执行顺序不同，那么比较就不可能进行，或者徒劳、无趣、极其困难——随你怎么说。更明确地说：只有当成对的计算能分解为按时间相继发生、彼此能够对应的动作，而相应程序文本也能同样分解为分别对应这些动作的指令时，比较这两个程序及它们可能引发的计算才有吸引力。

这是一个很强的条件。先举一个例子。

排除布尔条件检查的副作用，并假定“$B2$”的值不变（即执行“$S1$”或“$S2$”都不影响它），\textit{我们}可以确认以下两个程序等价：
\begin{equation}
	\label{eq:comparing-example-1}
	\begin{split}
	&\text{\hspace{-.4em}``\textbf{if }} B2\text{ \textbf{then}}\\
	&\quad \text{\textbf{begin while} } B1\text{ \textbf{do} } S1 \text{ \textbf{end}}\\
	&\text{\textbf{else}}\\
	&\quad \text{\textbf{begin while} } B1\text{ \textbf{do} } S2\text{ \textbf{end}''}
	\end{split}
\end{equation}

\noindent
以及
\begin{equation}
	\label{eq:comparing-example-2}
	\begin{split}
		&\text{\hspace{-.4em}``\textbf{while }} B1\text{ \textbf{do}}\\
		&\text{\textbf{begin if} } B2\text{ \textbf{then} } S1\text{ \textbf{else} } S2\text{ \textbf{end}''}
	\end{split}
\end{equation}

\label{pg:dilema}
第一个结构在初层上由选择子句控制执行顺序，第二个结构在初层上由重复子句控制执行顺序。我能够确认这些计算的输出等价，却不能在其他任何有用的意义上把它们看作等价。我不得不强迫自己接受这个结论：式（\ref{eq:comparing-example-1}）和式（\ref{eq:comparing-example-2}）“难以比较”。起初，这个结论让我非常恼火。后来，我逐渐把这种不可比性看作生活中的一个事实，也因此把它视为一个主要理由：在式（\ref{eq:comparing-example-1}）和式（\ref{eq:comparing-example-2}）之间选择，是一项相关的设计决定，不应未经仔细考虑就作出。正是它表面上的微不足道，使我留意到哪些因素应当影响这种选择。这些因素超出了本节范围，但我希望后面再谈。

再举一个稍微隐蔽些的不可比例子。

给定两个数组 $X[1:N]$ 和 $Y[1:N]$，以及布尔变量“equal”，要求编写程序，给“equal”赋予这样的布尔值：“两个数组逐元素相等”。空数组（即 $N = 0$）视为相等。

引入变量 $j$，并给“equal”赋予“前 $j$ 对元素中尚未发现差异”的含义，我们可以写出以下两个程序。
\begin{equation}
	\label{eq:comparing-example-3}
	\begin{split}
		&\text{\hspace{-.4em}``}j \coloneq 0;\text{ equal} \coloneq \text{\textbf{true}};\\
		&\text{\textbf{while} } j \neq N\text{ \textbf{do}}\\
		&\quad\text{\textbf{begin} } j \coloneq  j + 1;\text{ equal} \coloneq \text{equal \textbf{and} } (X[j] = Y[j]) \text{ \textbf{end}''}
	\end{split}
\end{equation}

\noindent
以及
\begin{equation}
	\label{eq:comparing-example-4}
	\begin{split}
		&\text{\hspace{-.4em}``}j \coloneq 0;\text{ equal} \coloneq \text{\textbf{true}};\\
		&\text{\textbf{while} } j \neq N\text{ \textbf{and} equal \textbf{do}}\\
		&\quad\text{\textbf{begin} } j \coloneq  j + 1;\text{ equal} \coloneq (X[j] = Y[j]) \text{ \textbf{end}''}
	\end{split}
\end{equation}

程序（\ref{eq:comparing-example-4}）与程序（\ref{eq:comparing-example-3}）的不同在于：一旦发现某一对元素不同，就终止重复。对于同一输入，两个程序的重复次数可能不同；因此，只有把程序最后两行视为描述一个整体动作、不再拆成子动作时，它们才在我们所说的意义上可比。但是，如果确实希望考虑两者都以重复结束这一事实，它们又有什么关系？为了弄清楚，我们来证明这两个程序的正确性。

对于数组 $X$ 和 $Y$，可以在 $0 \leqslant j \leqslant N$ 上定义如下 $N + 1$ 个函数 $\text{EQUAL}_j$：
\begin{equation}
	\label{eq:comparing-example-5}
	\begin{split}
		\text{当 } j = 0\qquad & \text{EQUAL}_j = \text{\textbf{true}} \\
		\text{当 } j > 0\qquad  & \text{EQUAL}_j = \text{EQUAL}_{j-1}\text{ \textbf{and} } (X[j] = Y[j]).
	\end{split}
\end{equation}

用这些函数表达，要求建立的净效果是
$$
\text{equal} = \text{EQUAL}_N.
$$

两个程序都维持关系
\begin{equation}
	\text{equal} = \text{EQUAL}_j
\end{equation}

\noindent
其中 $j$ 从 $0$ 开始递增。

我们很容易想把程序（\ref{eq:comparing-example-3}）和程序（\ref{eq:comparing-example-4}）看成同一个（抽象）程序（\ref{eq:comparing-example-7}）的两种细化：
\begin{equation}
	\label{eq:comparing-example-7}
	\begin{split}
		&\text{\hspace{-.4em}``}j \coloneq 0;\text{ equal} \coloneq \text{EQUAL}_0;\\
		&\text{\textbf{while} ``可能仍有：equal} \neq \text{EQUAL}_N\text{'' \textbf{do}}\\
		&\quad\text{\textbf{begin} } j \coloneq  j + 1;\text{ ``equal} \coloneq \text{EQUAL}_j\text{'' \textbf{end}''}
	\end{split}
\end{equation}

\noindent
其中，“可能仍有：equal $\neq$ EQUAL$_N$”代表某种尚未具体确定的基本操作。对它求值时，
$$
\text{equal} = \text{EQUAL}_j
$$

\noindent
成立；程序（\ref{eq:comparing-example-3}）和程序（\ref{eq:comparing-example-4}）的区别，在于它们用不同的判据保证“equal”取得最终值 EQUAL$_N$。

程序（\ref{eq:comparing-example-3}）中的判据很朴素，即
$$
j = N.
$$

\noindent
在被重复的语句开始时，
$$
\text{equal} = \text{EQUAL}_j
$$

\noindent
仍然成立。因此，执行“$j \coloneq j + 1$”之后，
$$
\text{equal} = \text{EQUAL}_{j-1}
$$

\noindent
成立；于是赋值语句
$$
\text{``equal} \coloneq \text{equal \textbf{and} } (X[j] = Y[j])\text{''}
$$

\noindent
就直接是递推关系（\ref{eq:comparing-example-5}）的转写。

要得到程序（\ref{eq:comparing-example-4}），则需要分析递推关系（\ref{eq:comparing-example-5}）。从中可以（再次用数学归纳法）推出，$\text{EQUAL}_j = \text{\textbf{false}}$ 蕴含 $\text{EQUAL}_N = \text{\textbf{false}}$，所以 $\text{EQUAL}_j = \text{\textbf{false}}$ 蕴含 $\text{EQUAL}_j = \text{EQUAL}_N$。如果出现这种情况，也能保证等式“$\text{equal} = \text{EQUAL}_N$”，由此便得到程序（\ref{eq:comparing-example-4}）。程序（\ref{eq:comparing-example-4}）中，被重复语句需要处理的（子）计算集合，限定为初始状态满足“$\text{equal} = \text{\textbf{true}}$”的那些计算。因此，在程序（\ref{eq:comparing-example-4}）中，赋值“$\text{equal} \coloneq \text{EQUAL}_j$”可以简写为
$$
\text{``equal} \coloneq (X[j] = Y[j])\text{''}
$$

现在就清楚了：为什么把式（\ref{eq:comparing-example-7}）引入为式（\ref{eq:comparing-example-3}）和式（\ref{eq:comparing-example-4}）的抽象，会造成误导。用“可能仍有：equal $\neq \text{EQUAL}_N$”，我们说明了某个布尔表达式为真和为假的含义，却没有给出表达式本身，这是很危险的。我们试图把式（\ref{eq:comparing-example-7}）理解成一个程序，在它自身的层次上，部分执行顺序尚未定义，而且随不同细化而变化。因此，我们试图把式（\ref{eq:comparing-example-7}）的最后几行同时看作式（\ref{eq:comparing-example-3}）和式（\ref{eq:comparing-example-4}）最后几行的模型；但这会误导，因为它们引发的计算无法建立一一对应关系。

关于我们认为不可比的程序，就谈到这里。可比程序的例子将在后面几节中出现。最后补充一句：我们说过，“成对的计算可以分解为按时间相继发生、彼此能够对应的动作”。但我们并未要求，配对动作的净效果必须相同！我们既可以比较完成同一工作的不同程序方案，也可以比较完成相似工作的不同程序。

\section[逐步构造程序的第一个例子]{逐步构造程序的第一个例子}
\label{sec:first-example-step-wise-program-composition}

在“\nameref{sec:undertanding-programs}”一节中，我强调了系统地安排执行顺序的必要性，使计算的结构能够体现在程序结构中；这样，我们就可以说，程序与计算共同形成结构。本节将尝试为“安排计算的结构”这个仍较模糊的概念，增添一些具体内容。这是第一次尝试：利用抽象能力，减轻枚举推理的负担；同时，它也将一以贯之地运用“\nameref{sec:undertanding-programs}”一节中提到的分解方式。

我不打算拿一个现成产品来展示所谓结构良好的程序，而要非常详细地描述这样一个程序的构造过程。这样做，是因为程序并不是现成存在的，恰恰相反，它们必须由人构造。我特别感兴趣的程序，是那些在我看来与我们的构造能力和理解能力相当适合的程序。

任务是指示计算机打印一张前一千个素数的表，规定 2 是第一个素数。

\begin{enumerate}[leftmargin=*, itemindent=2.2\parindent, label=\textit{注 }\arabic*.]
	\item 选择这个例子，是因为一方面，它足够困难，可以作为程序设计中某些问题的模型；另一方面，它的数学背景又简单、熟悉，不至于让问题本身占据我们的全部注意力。
	
	\item 我并不声称最终程序会是“最好的”，无论读者愿意用什么尺度来衡量。在本说明的先前版本中，余数是用除法计算的。至少两位读者对此反应激烈：“谁都知道，生成素数最有效的方法是埃拉托色尼筛法。”于是，他们再也读不下去了！
\end{enumerate}

我的基本做法，是以极小的步子构造程序，每次尽可能少作决定。问题分析逐渐深入，程序也随之进一步细化。

要构造一个算法，就必须用那些与我们已经充分理解的指令集合相对应的动作，组成所需要的计算。

程序最简单的形式是
\medskip


\noindent
描述 0：
\labeltext{描述 0}{tx:description-0}

\textbf{begin} “打印前一千个素数” \textbf{end}

\noindent
如果“打印前一千个素数”指的是那套已充分理解的指令集合中的一条指令，那么\ref{tx:description-0}就解决了问题。为了继续讨论，假定这条指令\textit{不在}该指令集合中。因此，我们必须构想一个由“更基本”的动作组成的计算，产生所期望的净效果。第一个方案是把素数生成与打印分开，于是提出\ref{tx:description-1}：
\medskip

\textbf{begin variable} “表 $p$”；
\labeltext{描述 1}{tx:description-1}

\quad “将前一千个素数填入表 $p$”；

\quad “打印表 $p$”

\textbf{end}，

\noindent
这表明，计算由两个按时间相继发生的动作组成，而且发生在只含一个变量的状态空间中，该变量叫作“表 $p$”。第一个动作给它赋值；第二个动作由它当时的值控制。

同样，如果“将前一千个素数填入表 $p$”与“打印表 $p$”都在那套已充分理解的指令集合中，而且“表 $p$”也是隐含可用的资源，那么问题就解决了。为继续讨论，我们再次假定并非如此。这意味着，下一步细化时，必须说明如何通过另外两个（子）计算产生这两个动作的效果。此外，还必须决定：怎样表示“表 $p$”的中间值所包含的信息——这个对象目前还相当不明确。

在继续之前，我想强调：写下描述 1 时，我们作出的决定有多么少，考虑到原始问题说明的内容又有多么少。我们假定某种形式的“表 $p$”资源可用，因此能在开始打印前，算出前一千个素数。在这个假定下，我们利用了一个事实：素数的计算可以独立于打印来构思。对于原始问题说明，除此以外，我们考虑的内容大体只限于：至少存在一千个不同的素数（要使问题说明有意义，就必须假定这一点）。在这个阶段，“素数”这个概念究竟是什么意思，仍然不太重要。对于打印输出的具体版式要求，我们也完全没有作出承诺。显然，这种做法的优点在于：原始问题中这两个相当独立的方面所带来的后果，似乎已经分别归入两个组成动作各自的细化中。这表明，在应用“分而治之”这一黄金原则时，我们多少取得了成功。

不过，继续讨论时，我们必须问自己：现在，这两个子计算究竟能在多大程度上彼此独立地构思？更精确地说：“我们是否已经达到这样的阶段：可以让两位程序员互不依赖地设计两个子算法，以引发相应的子计算？”

既然这两个动作已不能视为由充分理解的指令所调用，那么变量“表 $p$”也就不能再视为隐含可用的资源。正如必须把动作分解为子动作，我们也必须选择怎样组成变量“表 $p$”：选择什么数据结构，来表示从第一个动作经由“表 $p$”传递给第二个动作的信息。这个决定总得在某个时刻作出，问题在于：“什么时候？”以及“怎样作出？”

原则上，似乎有两条路。第一条是尽量推迟决定如何把“表 $p$”组织成一些更中性、较少依赖特定问题的部件。如果推迟对“表 $p$”结构的决定，下一步就要细化一个或两个动作。为此，可以假定在这个仍然神秘的对象\mbox{“表 $p$”}上存在一套合适的操作；最后，再汇总这些操作，根据它们的要求，设计最合适的“表 $p$”结构。

另一条路，是尝试就在此时此地决定“表 $p$”的结构。一旦前一千个素数表的表示方式确定下来，两个动作的细化就能相当独立地进行。

两条路都棘手。例如，第一个子计算的算法是否合适，很大程度上取决于所假定的“表 $p$”操作能否容易、优雅地实现；如果其中一个或几个操作笨拙到无法接受，整座大厦就会倒塌。反过来，如果过早决定“表 $p$”的结构，也很可能发现子计算随后变得别扭。这个问题无法绕开：在优雅的程序中，“表 $p$”的结构与引用它的计算必须相互匹配。我认为，高效程序员的行为可以这样描述：努力先作最容易的决定，即用尽可能少的探究——试错、反复相互调整等——获得尽可能充分的理由，相信自己将来不会后悔这个决定。

为了不让这番讨论过分冗长，我们假定程序员鼓起勇气，决定现在首先确定“表 $p$”的结构。一旦采取这个立场，立即会出现两种方案。一方面，我们可以利用这样的事实：“前 1000 个素数的表”不只是一张含有一千个数的表——不像工厂里 1000 名员工的月工资表——它的所有数还互不相同。利用这一点，可以用线性布尔数组组织信息，将连续的数组元素对应于连续的自然数，指示相应自然数是否为素数。数论给出第一千个素数数量级的估计，由此可以界定一个足够长的数组长度。如果这样组织数据，就准备了一个容易回答“$n$（小于最大值）是否为素数？”的机制。另一方面，我们可以选择整数数组，依次列出素数。（如果需要\textit{事先}给定整数数组元素的最大值，这里也会用到数论得到的同一个估计。）后一种形式，则建立了适合回答“第 $k$ 个素数的值是什么，其中 $k \leqslant 1000$？”的机制。

我们让程序员勇敢地选择后一种表示。这对打印操作似乎很合适，因为要求的是打印素数，而不是打印自然数并标记它们是不是素数。对计算阶段，它也似乎很合适；这里要给程序员一点预见：分析一个自然数是否为素数，会与能否找到这个数的素因子有关。

于是，程序细化的下一阶段，是仔细说明这个仍然神秘的对象“表 $p$”采用何种表示约定，并按这一约定重新定义两个操作。

约定如下：“表 $p$”中的信息，由“\textbf{integer array} $p[1 : 1000]$”各元素的值表示；当素数按数值递增顺序排列时，对 $1 \leqslant k \leqslant 1000$，$p[k]$ 等于第 $k$ 个素数。（如果整数的最大值已隐含给定，我们假定数论能让我们断定它足够大。）

现在，要描述这一新的细化，就遇到一个新困难。\ref{tx:description-1}之所以能写成单个程序，是因为它细化的是\ref{tx:description-0}中的单个已命名动作“打印前一千个素数”。（用更常见的说法，\ref{tx:description-1}可以写成一个过程体。）但下一层不再如此：我们必须在某种意义上同时细化三个已命名的实体——“表 $p$”和两个动作。因此，需要想出一套标识用语，说明什么细化了什么。

为继续讨论，我们暂时作一个试探性的建议：\ref{tx:description-0}是一段有效文本，用单个已命名动作“打印前一千个素数”来表达；用代码 $0a$ 标识这个动作。

\hyperref[tx:description-1]{描述 1}之所以称为“\hyperref[tx:description-1]{1}”，是因为它是\ref{tx:description-0}的下一层细化。它细化了 $0a$\textemdash{}\ref{tx:description-0}唯一使用的项\textemdash{}自身则用三个已命名实体表达；我们给它们附上代码：

\noindent
“表 $p$”\hfill 1a

\noindent
“将前一千个素数填入表 $p$”\hfill 1b

\noindent
“打印表 $p$”\hfill 1c

\noindent
代码以 1 开头，因为\ref{tx:description-1}是用它们表达的；“a”“b”“c”则用于区分。

现在必须描述“表 $p$”内信息所采用的表示约定，而这个约定涉及 1a、1b、1c 三个元素。因此，把它叫作\ref{tx:description-2}；其中应分别描述三个元素（用等号作分隔符）：
\medskip

\noindent
描述 2：
\labeltext{描述 2}{tx:description-2}
\nopagebreak

\noindent
1a $=$ “\textbf{integer array} $p[1 : 1000]$”

\noindent
1b $=$ “对 $k$ 从 1 到 $1000$，使 $p[k]$ 等于第 $k$ 个素数”

\noindent
1c $=$ “对 $k$ 从 1 到 1000，打印 $p[k]$”。

\hyperref[tx:description-2]{描述 2}用三个已命名实体表达，按显而易见的次序，赋予它们代码 2a、2b、2c。（就代码而言，\ref{tx:description-2}内容很少：它只是说，对 1a、1b、1c，我们分别选择了细化 2a、2b、2c。）

\noindent
\textit{说明}。在表示“表 $p$”内的信息时，我们没有利用每个待打印的值只出现一次这个事实，也没有利用它们按数值递增顺序出现这个事实。反过来，这意味着：名为 2c 的动作，被看作“打印任意一组一千个整数值”的一个具体实例（它也可以是一千名编号员工的月工资表！）。这个例子中，打印动作的净效果，与前一千个素数一样，是唯一确定的；但我们把它构想成某个更大事件类别中的具体实例。进一步细化 2c 时，处理的是整个类别；这一类中的具体实例，则由数组 $p$ 各元素的值确定。当人们说“定义接口”时，我常觉得，他们忽略了其中预设的推广，忽略了对一类“可能”动作的构想。

如果 2b 和 2c 都在那套已充分理解的指令集合中（于是 2a 也是隐含可用的资源），整个问题就解决了。为继续讨论，再次假定并非如此。于是，我们需要为动作 2b 和 2c 构想子计算。但现在，由于引入了第 2 层，2b 与 2c 各自的细化可以独立设计。

\emph{2b 的细化：“对 $k$ 从 $1$ 到 $1000$，使 $p[k]$ 等于第 $k$ \emph{个}素数”}。

我们要找的是描述 2b1，即 2b 的第一次细化。为了表明 2b 与 2c 的细化彼此独立，我们在 2b 后另起编号，而不把下一个描述叫作“3 什么”。

在描述 2b1 中，要给出一个算法，描述数组 $p$ 各元素怎样得到自己的值。例如，这意味着必须说明赋值的顺序。第一次细化时，我们就只说明这一点，最好不说明其他事情。一个显而易见却荒唐的版本，开头如下（“版本号”放在括号里）：

\noindent
2b1(l)：
\nopagebreak

\noindent
\textbf{begin} $p[1] \coloneq 2$; $p[2] \coloneq 3$; $p[3] \coloneq 5$; $p[4] \coloneq 7$; $p[5] \coloneq 11$;\ \dots\ \textbf{end}

\noindent
这意味着程序员已经知道前一千个素数的表。我们不再考虑这个版本，因为这就意味着程序员几乎根本不需要机器。

第一个素数已给定 $(=2)$，并假定程序员不知道第一千个素数。于是，填充数组 $p$ 各元素最自然的顺序，似乎就是下标递增的顺序。如果只表达这一点，例如可以得到

\noindent
2b1(2)：

\noindent
\textbf{begin integer} $k$, $j$; $k \coloneq 0$; $j \coloneq 1$;

\noindent
\quad \textbf{while} $k < 1000$ \textbf{do begin} “增加 $j$，直到下一个素数”；

\noindent
\quad\quad $k \coloneq k + 1$; $p[k] \coloneq j$ \textbf{end}

\noindent
\textbf{end}

把 $k$ 理解为已找到的素数个数，并核实第一个素数 $(=2)$ 确实是大于 1（$=$ $j$ 的初始值）的最小素数，就很容易用数学归纳法证明 2b1(2) 正确（假定存在足够多的素数）。

如果“增加 $j$，直到下一个素数”所描述的操作——称为 2b1(2)a——在指令集合中，那么描述 2b1(2) 就是一个完善的程序。但假定它不在。这样，下一次细化时，就必须表达怎样增加 $j$（同样，最好不表达其他事情）。于是得到第 2b2(2) 层的描述：

\noindent
2b1(2)a $=$

\noindent
\textbf{begin boolean} $j$prime;

\noindent
\quad \textbf{repeat} $j \coloneq j + 1$;

\noindent
\quad\quad “赋予 $j$prime 这样的含义：$j$ 是素数”

\noindent
\quad \textbf{until} $j$prime

\noindent
\textbf{end}

\noindent
\textit{说明}。这里使用 repeat-until 子句，是为了表明 $j$ 总要至少增加一次。

其正确性同样几乎无可怀疑。不过，如果假定程序员知道除 2 以外的所有素数都是奇数，那么由于效率低，他大概不会满意这个版本。为这种“缺乏预见”所付出的代价，就是修改版本 2b1(2)。单独处理素数 2 之后，循环就只需处理奇素数。替代 2b1(2)，我们得到

\noindent
2b1(3)：

\noindent
\textbf{begin integer} $k$, $j$; $p[1] \coloneq 2$; $k \coloneq 1$; $j \coloneq 1$;

\noindent
\quad\textbf{while} $k < 1000$ \textbf{do}

\noindent
\quad\quad\textbf{begin} “增加奇数 $j$，直到下一个奇素数”；

\noindent
\quad\quad\quad $k \coloneq k + 1; p[k]: = j$

\noindent
\quad\quad\textbf{end}

\noindent
\textbf{end}

\noindent
对引号中的操作——姑且称为“2b1(3)a”——作类似细化，得到第 2b2(3) 层的描述：

\noindent
2b1(3)a $=$

\noindent
\textbf{begin boolean} $j$prime;

\noindent
\quad\textbf{repeat} $j \coloneq j + 2;$

\noindent
\quad\quad“对于奇数 $j$，赋予 $j$prime 这样的含义：$j$ 是素数”；

\noindent
\quad\textbf{until} $j$prime

\noindent
\textbf{end}

刚才在两个描述层次间来回调整，实质上不过是为了方便，调整总体结构与必须适合这个结构的基本操作之间的接口。这种往复调整、这种试错方式，确实不理想；但当我们缺乏预见、又不得不依次作决定时，我看不到别的办法。可以把这些努力看作一些试验：以相当低的代价，探索接口大概选在哪里最方便。

\noindent
\textit{说明}。2b1(2) 和 2b1(3) 都可以粗略描述为

\noindent
\textbf{begin} “将表 $p$ 和 $j$ 设为初始值”；

\noindent
\quad\textbf{while} “表 $p$ 尚未填满” \textbf{do}

\noindent
\quad\textbf{begin} “增加 $j$，直到下一个待加入的素数”；

\noindent
\quad\quad “将 $j$ 加入表 $p$”

\noindent
\quad\textbf{end}

\noindent
\textbf{end}

\noindent
但我们不这样做，因为两个版本的执行顺序不同（参见“\nameref{sec:comparing-programs}”），我们把它们看作“不可比”。选择 2b1(3)，就意味着认定试用版本 2b1(2) 和 2b1(1) 一样，不再适用，因此舍弃。

从 2b1(2) 改为 2b1(3)，理由是更深的细化层次能获得效率收益。这项收益在第 2b2 层获得，因为现在 $j$ 可以每次加 2；在 2b2(3) 层尚未确定的基本操作中，这项收益也会体现出来：用于“对于奇数 $j$，赋予 $j$prime 这样的含义：$j$ 是素数”的算法，只需分析 $j$ 的奇数值。

再说一次：在 2b2(3) 中，我们用一个算法细化了 2b1(3)。只要“对于奇数 $j$，赋予 $j$prime 这样的含义：$j$ 是素数”——称为“2b2(3)a”——在那套已充分理解的指令集合中，这个算法就解决了问题。现在假定它不在；换句话说，需要引发一个计算，判断给定的奇数 $j$ 是否有因子。直到这一步，代数才真正进入讨论。这里，我们利用已知事实：只需尝试素因子；而且，待试的素数都已经能在数组 $p$ 已填充的部分中找到。

我们使用以下事实：

\begin{enumerate}[wide, nosep, label=(\arabic*)]
	\item $j$ 是奇数，因此应尝试的最小潜在因子为 $p[2]$，即大于 2 的最小素数。

	\item 如果 $p[\text{ord}]$ 是平方大于 $j$ 的最小素数，那么应尝试的最大素数是 $p[\text{ord} - 1]$。
\end{enumerate}

（这里，我还用了一个事实：平方大于 $j$ 的最小素数已经可以在表 $p$ 中找到。对这个程序的一个早期版本，我把这一点视为当然；以下怀着谦逊，引用 Don Knuth 当时的评论：

\begin{addmargin}[\parindent]{0em}
“这里你犯了一个严重的遗漏！你的程序使用了一个深刻的数论结果：如果 $p_n$ 表示第 $n$ 个素数，那么\textit{总有}
\end{addmargin}
$$
p_{n+1}<p_n^2.\text{”}
$$

\noindent
我认错。）

如果这个集合非空，我们就有机会找到因子；一旦找到因子，就可以停止考察这个 $j$ 值。必须决定按什么顺序尝试集合中的素数；我们选数值递增顺序，因为素数越小，成为 $j$ 因子的概率越大。

当 ord 的值已知时，可以为“对于奇数 $j$，赋予 $j$prime 这样的含义：$j$ 是素数”给出第 2b3(3) 层的描述：

\noindent
2b2(3)a =

\noindent
\textbf{begin integer} $n$; $n \coloneq 2$; $j$prime $\coloneq$ \textbf{true};

\noindent
\quad\textbf{while} $n <$ ord \textbf{and} $j$prime \textbf{do}

\noindent
\quad\textbf{begin} “赋予 $j$prime 这样的含义：$p[n]$ 不是 $j$ 的因子”；$n \coloneq n + 1$

\noindent
\quad\textbf{end}

\noindent
\textbf{end}

但这个版本假定，作为 $j$ 的函数，ord 的值已经知道。我们本来也可以这样开始细化：

\noindent
\textbf{begin integer} $n$, ord;

\noindent
\quad ord $\coloneq 1$; \textbf{while} $p[\text{ord}] \uparrow 2 \leqslant j$ \textbf{do} ord $\coloneq$ ord $+ 1$;

\noindent
\quad \dots

\noindent
即每次需要“ord”的值，都重新计算。这里似乎适合用存储空间换取计算时间：与其每次需要这个函数就重新计算，不如引入一个额外变量 ord，保存它的当前值。设置 $j$ 时要设置它，改变 $j$ 时要调整它。

可惜，这迫使我们重新编写一些程序。一种办法是：在引入 $j$ 的同时引入整数变量 ord，然后扫描程序，在每个操作 $j$ 的地方插入适当的 ord 操作。我不喜欢这种做法，因为在引入 $j$、并赋予它含义的那一层，函数“ord”无关紧要。因此，我们将努力只在合适的层次引入 ord，并且非常谨慎。

对于 2b：“对 $k$ 从 1 到 1000，使 $p[k]$ 等于第 $k$ 个素数”，我们写下（与第 2b1(3) 层类似）

\noindent
第 2b1(4) 层：

\noindent
\textbf{begin integer} $k$, $j$; $p[1] \coloneq 2$; $k \coloneq 1$;

\noindent
\quad “将 $j$ 置为一”；

\noindent
\quad\textbf{while} $k < 1000$ \textbf{do}

\noindent
\quad\textbf{begin} “增加奇数 $j$，直到下一个奇素数”；

\noindent
\quad\quad $k \coloneq k + 1$; $p[k] \coloneq j$

\noindent
\quad\textbf{end}

\noindent
\textbf{end}

\noindent
它用以下实体表达：

\noindent
2b1(4)a\tabto{5em} “增加奇数 $j$，直到下一个奇素数”

\noindent
2b1(4)b\tabto{5em} “将 $j$ 置为一”。

在下一层，只引入 2b1(4)a 的子计算；另一个则向下传递。

\noindent
第 2b2(4) 层：

\noindent
2b1(4)a $=$

\noindent
\textbf{begin boolean} $j$prime;

\noindent
\quad\textbf{repeat} “将 $j$ 加二”；

\noindent
\quad\quad “对于奇数 $j$，赋予 $j$prime 这样的含义：$j$ 是素数”

\noindent
\quad\textbf{until} $j$prime

\noindent
\textbf{end;}

\noindent
2b1(4)b $=$ 2b2(4)b

\noindent
它用以下实体表达：

\noindent
2b2(4)b\tabto{5em} 仍表示“将 $j$ 置为一”

\noindent
2b2(4)c\tabto{5em} “将 $j$ 加二”

\noindent
2b2(4)d\tabto{5em} “对于奇数 $j$，赋予 $j$prime 这样的含义：$j$ 是素数”。

直到下一层，我们才需要谈 ord。因此，现在写下

\noindent
第 2b3(4) 层：\textbf{integer} ord;

\noindent
2b2(4)b =

\noindent
\textbf{begin} $j \coloneq 1$; “设置 ord 的初始值” \textbf{end};

\noindent
2b2(4)c =

\noindent
\textbf{begin} $j \coloneq j + 2$; “调整 ord” \textbf{end};

\noindent
2b2(4)d =

\noindent
\textbf{begin integer} $n$; $n \coloneq 2$; $j$prime $\coloneq$ \textbf{true};

\noindent
\quad\textbf{while} $n <$ ord \textbf{and} $j$prime \textbf{do}

\noindent
\quad\quad\textbf{begin} “赋予 $j$prime 这样的含义：$p[n]$ 不是 $j$ 的因子”；

\noindent
\quad\quad\quad $n \coloneq n + 1$

\noindent
\quad\textbf{end}

\noindent
\textbf{end}

\noindent
它用以下实体表达：

\noindent
2b3(4)a\tabto{5em} “设置 ord 的初始值”

\noindent
2b3(4)b\tabto{5em} “调整 ord”

\noindent
2b3(4)c\tabto{5em} “赋予 $j$prime 这样的含义：$p[n]$ 不是 $j$ 的因子”。

下一层给出两个彼此独立的细化。（注：本来也可以把它们放在先后相继的层次中，但那样就必须给这两个层次施加一个任意顺序。也可以尝试分别处理这两个细化，就像分别处理 2b 和 2c 那样；但我们觉得，这个明确分开的决定现在作出还稍早。）我们将表达：

\begin{enumerate}[leftmargin=2\parindent, label=(\arabic*)]
	\item ord 是 $j$ 的非递减函数，而 $j$ 的值只会增加，因此调整 ord 意味着有条件地增加它；

	\item $p[n]$ 是否为 $j$ 的因子，由余数是否等于零来确定。
\end{enumerate}

\noindent
由此得到

\noindent
第 2b4(4) 层：

\noindent
2b3(4)a $=$ 2b4(4)a

\noindent
2b3(4)b $=$

\noindent
\textbf{begin while} “ord 太小” \textbf{do} “将 ord 加一” \textbf{end};

\noindent
2b3(4)c =

\noindent
\textbf{begin integer} $r$;

\noindent
\quad “使 $r$ 等于 $j$ 除以 $p[n]$ 的余数”；

\noindent
\quad $j$prime $\coloneq (r \neq 0)$

\noindent
\textbf{end}

\noindent
它用以下实体表达：

\noindent
2b4(4)a\tabto{5em} 仍表示“设置 ord 的初始值”

\noindent
2b4(4)b\tabto{5em} “ord 太小”

\noindent
2b4(4)c\tabto{5em} “将 ord 加一”

\noindent
2b4(4)d\tabto{5em} “使 $r$ 等于 $j$ 除以 $p[n]$ 的余数”

如果有内置除法，可以假定实现“使 $r$ 等于 $j$ 除以 $p[n]$ 的余数”很容易。至于独立处理 2b4(4)d 细化的情况，就留给有兴趣的读者。为使算法有一个出人意料的转折，我们假定没有方便的求余计算。此时，算法
$$
\text{``}r: = j;\text{ \textbf{while} } r > 0\text{ \textbf{do} } r \coloneq r - p[n]\text{''}
$$

\noindent
会得到（非正的）余数，但从计算时间来看，这很不合适。因此又需要引入一些额外的表存数据，方式类似于引入“ord”时的做法。

我们想知道，对于 $n < \text{ord}$，给定的 $j$ 值是否为 $p[n]$ 的倍数。为帮助分析，引入第二个数组；其元素可存储各个相继素数的倍数，尽可能方便地靠近 $j$。为了给出数组大小，希望知道 ord 的上界；当然，1000 是安全的，但数论给出了 30 这个安全上界。因此引入

\noindent
\textbf{integer array} mult$[1 : 30]$

\noindent
并作如下约定：对于 $n < \text{ord}$，mult$[n]$ 是 $p[n]$ 的倍数，且满足关系
$$
\text{mult}[n] < j + p[n]
$$

\noindent
增加 $j$ 时，这一关系始终保持成立。每次想检查 $p[n]$ 是否为 $j$ 的因子，就只要
$$
\text{mult}[n] < j.
$$

\noindent
还成立，就反复给 mult$[n]$ 加上 $p[n]$。增加完成后，mult$[n] = i$ 是 $j$ 为 $p[n]$ 倍数的充分必要条件。

ord 的最大值很小，还有另一个后果：“ord 太小”这一检查可以表达为
$$
\text{``}p[\text{ord}] \uparrow 2 \leqslant j\text{''}
$$

\noindent
但对于同一个 ord 值，这项检查需要执行许多次。可以假定，引入一个值等于 $p[ord] \uparrow 2$ 的变量（称为“square”），就能提高速度。

于是，得到最终的

\noindent
第 2b5(4) 层：

\noindent
\textbf{integer} square; \textbf{integer array} mult$[1 : 30]$;

\noindent
2b4(4)a $=$

\noindent
\textbf{begin} ord $\coloneq 1$; square $\coloneq 4$ \textbf{end};

\noindent
2b4(4)b =

\noindent
\quad (square $\leqslant j$);

\noindent
2b4(4)c =

\noindent
\textbf{begin} mult$[\text{ord}] \coloneq$ square; $\text{ord} \coloneq \text{ord} + 1$; square $\coloneq p[\text{ord}] \uparrow 2$ \textbf{end};

\noindent
2b4(4)d =

\noindent
\textbf{begin while} mult$[n] < j$ \textbf{do} mult$[n]\coloneq$ mult$[n] + p[n]$; $r\coloneq j -$ mult$[n]$ \textbf{end}

\noindent
这使我们的计算已经接近埃拉托色尼筛法的一种实现！

\noindent
\textit{注}。在 2b4(4)d 的细化中，把 mult$[n]$ 与 $j$ 的当前值比较时，mult$[n]$ 会尽可能增加到满足要求。这本来可以用 $2\times p[n]$ 为步长，因为我们只提交奇数 $j$，所以只对 $p[n]$ 的奇数倍感兴趣。（mult$[1]$ 的值一旦设置，就一直等于 4。）

2c“对 $k$ 从 1 到 1000，打印 $p[k]$”的细化留给读者。我建议把表打印成五页，每页四栏，每栏包含五十个连续的素数。

\vspace*{-1em}
\begin{center}
	$\star$\quad\quad  $\star$\\
	\nopagebreak
	
	\vspace*{-.5em}
	$\star$
\end{center}
\vspace*{-1em}

到这里，我完成了本节开始时宣布要做的事：“非常详细地描述这样一个〔结构良好的〕程序的构造过程”。最后再说几点意见。

最引人注意的是：我们对一个非常简单的程序，竟然讨论了如此长的篇幅。即便考虑到实质上做了两件事——构造程序，以及详细讨论促成它的各种考虑——这个长度也确实超出了我的喜好和愿望。真正让我困扰的，与其说是后一部分的长度（作家可以写整部小说来描述人的行为），不如说是各个层次的文本长度。因此，可以预见，记法技术将成为我们的主要关注之一。

但也有令人鼓舞的经验。我们充分承认，可怜的程序员不可能一下子决定所有事情；在这个前提下，我们相当成功地做到，每次作一个决定，逐步构造程序。例子中，当一些重大决定仍悬而未决时，不少程序部分就已经写成最终形式：无论最后那些决定采用哪种方案，较早层次的代码仍然有效。从程序可管理性的要求来看，这很令人鼓舞。

\section[论程序族]{论程序族}

上一节中，我们考虑了如何为给定任务设计程序，但把最终程序看成了一个孤立对象：一个独立存在、只按自身优点评价的结构。这个结构来自连续的分解；结构的目的，是让程序的正确性能够得到证明，而不需要过度的智力劳动。

本节将说明：为什么与其把程序看成孤立对象，我更愿意把它看成“相关程序”族的一员。用传统术语来说，相关程序可以是完成同一任务的不同方案，也可以是完成相似任务的相似程序。

程序员为什么不能只关注自己要写的那个程序，而还要考虑整个程序族？首先，除非能够把自己的行动说成是在一组同样可能的行动中作出的有意选择，否则很难声称自己知道在做什么。但如果要充分承认构造大型复杂程序的特殊困难，这还有一个非常实际的理由。（而且，我们\textit{必须}承认这些特殊困难：经验表明，即使某人在某种规模上已证明自己能做得出色，也绝不能保证他遇到更大规模的工作时，不会弄得一团糟。）

大型程序的一个特点，当然是它们在使用期间必须修改。一个很常见的原因是：程序虽然逻辑正确，引发的计算却不令人满意，例如在一个或几个数量指标上不合要求。第二个原因是：程序虽然逻辑正确，甚至很好地满足了原始要求，却是对一个不太对的问题给出的完美解答；于是必须重新表述问题，并相应调整程序。

一种天真的做法，是认为我们必须能修改现有程序（奇怪的说法“程序维护”已经为此确立下来）。这样，任务就被看成文本操作。顺便回想一下，这项需要甚至曾被用作一个论据，支持用穿孔卡片、而不是纸带来输入程序文本。然而，实际修改程序文本只是事务性工作，可以用许多不同方式完成。我要说的是：如果我们主要把程序文本当作线性的符号序列来把握，那么随着文本越来越长，确定并描述哪些内容需要修改，往往会难到无法完成。

如果一个程序必须有两个不同版本，我宁愿不要把其中一个的（文本）看成另一个的（文本）的修改结果。更理想的是：在某种意义上，把它们看成同一个祖先的不同后代；这个祖先代表一个多少有些抽象的程序，体现两个版本的共同部分。希望能在（前期）文档中，轻易认出这个共同祖先。我们的目标是

\begin{enumerate}[wide, nosep, label=(\arabic*)]
	\item 两个版本尽可能共享各自的正确性证明；

	\item 两个版本尽可能以机械方式共享共同（或“相同”）代码；

	\item 受修改影响的区域已经很好地隔离。如果版本转换需要靠人脑决定、并在文本各处零散修改，就不满足这个条件。
\end{enumerate}

这是一个很高的目标。它来自程序修改与程序构造两项任务之间可能存在的相似性：当程序逐步构造到某个中间细化阶段时，已经写下的内容，实际上就是所有进一步细化所得程序的一个合适“共同祖先”。“需要改变的决定”和“尚未作出的决定”是相似的：两种情况下，我们面对的都是把这个决定抽象掉之后所剩下的部分。

我们的经验还提供了第二个启发来源。在由外向内、逐步细化的程序构造过程中，早期不仅推迟决定某些事情怎样做，也推迟承诺究竟要做什么。随着细化推进，实际问题说明中的更多细节才逐渐纳入考虑。（后面的例子会比素数表问题更清楚地展示这一点。）因此，前几个细化层次同样适用于一整类问题说明。

换句话说，逐步构造的方法建议：即使任务已经定义清楚，开始时仍应忽略问题说明的某些方面。这意味着，程序员不把给定任务视为一件孤立的待办事项，而是把它看作整个任务族的一员，并对给定问题作适当推广。然后陆续加入更多细节，最终把算法确定为给定问题的解。

这些都是众所周知的，每位称职的程序员一直在这样做。但我仍出于几种原因强调它。第一，如果问题说明十分繁复，无法一眼把握，程序员就\textit{必须}这样接近并剖析它（参见“\nameref{sec:our-inability-to-do-much}”一节）。第二，即使给定问题已完全定义清楚，也应尽可能预见问题说明将来的变化，并在能力所及的范围内作出安排；这是一项明智的预防措施。这并不是让人把程序写得如此“通用”，以至于效率低到无法接受；对问题说明作出考虑不周的推广，就容易出现这种情况（如果推广由销售部门指定，更容易如此！）。但根据我的经验，即使在传统程序设计中，寻找那些可行、又可能有用的推广，也很值得做：这种思考可以清楚地指引最终程序应采用怎样的结构。而这样的思考，说到底就是……多少明确地构想出整个程序族！

在前面“\nameref{sec:realiability-of-mechanisms}”一节中，我们从必须能够证明程序正确性这一要求，提出了精心安排程序结构的必要性。本节又遇到另一个理由：程序结构应当预先照顾到它的适应和修改。程序不仅应通过结构体现我们对它的理解，还应让人从结构中清楚看出，哪些调整能够顺利完成。谢天谢地，这两项要求是相辅相成的。

\section[论用存储空间换取计算速度]{论用存储空间换取计算速度}
\label{sec:trading-storage-for-computation-speed}

在今天的顺序计算机中（1969 年春），可以区分两个主要部件：主动部件（处理器）和被动部件（存储器）。主动部件的特定职责是快，被动部件的特定职责是大。下面的讨论假定，这种功能分工会保持足够长的时间，以至于研究其后果仍有意义。

从程序员的角度看，存储空间和计算时间是两种不同资源。我认为，分配这些资源——也就是在两者之间分担负荷——是程序员的责任之一，而非系统的责任。本节就讨论这项责任的后果。这里不讨论估计各项负荷的技术，即不提供影响程序员选择的数量标准；我们讨论的是：程序员可选的不同方案之间，有怎样的逻辑关系。

\noindent
\textit{注}。把某些选择交给系统，并非不可想象。不过，除了最简单的情况，设计方案并确认其等价性，似乎都需要程序员作出数学上的创造。本篇不讨论将这种问题求解活动自动化的任何尝试。

最简单的情况是：某个计算经常需要“FUN(arg)”的值，其中“FUN”是给定的可计算函数，定义在一个或多个已存储变量的当前值上；这些变量统称为“arg”。程序版本 $A$ 只存储 arg 的值，每次需要 FUN(arg) 时再计算它。版本 $B$ 则引入额外变量，比如“fun”，其唯一目的，是记录与 arg 当前值对应的“FUN(arg)”值。

版本 $A$ 中写
\medskip

{
	\setlength{\parindent}{8em}
	\hspace{-.4em}``arg $\coloneq \dots$''（即对 arg 赋值）
	
}
\medskip

\noindent
的地方，版本 $B$ 写
\medskip

{
	\setlength{\parindent}{8em}
	\hspace{-.4em}``arg $\coloneq \dots$; fun $\coloneq$ FUN(arg)''
	
}
\medskip

\noindent
从而维持关系
\medskip

{
	\setlength{\parindent}{8em}
	fun = FUN(arg)
}
\medskip

\noindent
由于这个关系成立，版本 $A$ 中每个需要对 FUN(arg) 求值的地方，版本 $B$ 都改为读取变量 fun 的当前值。

偏好版本 $B$ 而非版本 $A$，可能有两个理由。如果请求 FUN(arg) 值的频率，比对 arg 赋值的频率更高，那么版本 $B$ 可能需要较少计算时间。必要时，还可以引入另一个（布尔）变量“fun up to date”，表示是否假定关系“fun = FUN(arg)”成立，以进一步改进这一方法。于是，对 arg 赋值时，一并执行
\medskip

{
	\setlength{\parindent}{8em}
	\hspace{-.4em}``fun up to date $\coloneq$ \textbf{false}'';
}
\medskip

\noindent
每次需要 FUN(arg) 的值时，检查这个布尔变量，就知道是否必须重新计算 FUN(arg)。如果必须，就把算出的值赋给 fun，并按“fun up to date”的含义将它设为 \textbf{true}。把最后这个程序版本叫作 $C$。这三个程序仅在版本 $A$ 对 arg 赋值或使用 FUN(arg) 值的地方有所不同；显然，就输出而言，它们等价。通过机械方式，从版本 $A$ 得到版本 $B$ 或 $C$，也并非不可想象。

不过，情况常常没有这么简单。现在来谈引入“fun”变量的第二个理由。对于任意 arg 值，从头计算 FUN(arg) 往往很不合算；而当 arg 的值改变时，计算 FUN(arg) 的值怎样改变，却容易得多。在这种情况下，调整“fun”的值与函数依赖本身的性质，有着更紧密的联系，而不只是下面写法所暗示的关系：
\medskip

{
	\setlength{\parindent}{8em}
	\hspace{-.4em}``arg $\coloneq\dots$; fun $\coloneq$ FUN(arg)''
}
\medskip

这种可能性，往往不仅与函数依赖的性质紧密相关，还与计算进行过程中“变量 arg 的历史”有关！在素数表程序中，我们已经看到一个非常鲜明的例子（参见“\nameref{sec:first-example-step-wise-program-composition}”）：引入“ord”，它在函数意义上依赖于“$j$”；具体说，“ord”是满足
\medskip

{
	\setlength{\parindent}{8em}
	$p[\text{ord}]$ $\uparrow$ $2 > j$
}
\medskip

\noindent
的最小值。由于“$j$”随时间单调增加，调整“ord”就成了一项很方便的操作。

在理解程序时，我希望这些存储冗余信息的额外变量，能够明确地被识别为这样的变量，即便相关函数关系有些不完全确定，像同一例子中的“mult”表那样。我很倾向于把这些程序看成较抽象程序的优化细化，即使为了达到现实可用的性能，这些额外变量所带来的优化必不可少。从效率角度看，某个额外变量可能如此重要，以至于构想一个把它的存在抽象掉的层次，会被人觉得是不负责任的白日梦。人们常常把操作这类额外变量的方式，看作算法的主体；我们的数学创造力，也常常在那里结出成果。但我要强调的是：虽然要得到性能现实可用的程序，至少存在一种这样的优化细化很重要，可进一步检查后，常常会发现，这种优化细化即使在最粗略的层次上，也远不唯一。

\noindent
\textit{注}。我记得有一个程序，它的额外信息冗余到了这样的程度：不仅能从“arg”的值推出“fun”的值，反过来也可以。忽然，“fun”和“arg”的关系变成了对称关系，而一个问题令我十分困扰：我凭什么如此不对称地处理它们？那个程序生成一道组合谜题的全部解。进一步检查发现，还有第二道组合谜题，而且可以证明，两道问题的解之间存在一一对应关系。如果我求解的是第二道组合问题，就会发现“fun”和“arg”的角色互换了！在传统程序设计中，这种函数依赖并不显式展示，两道谜题很可能用完全相同的程序求解；而我构造了两个结构不同的程序。我认为这样才对，因为同一个程序用来解两道谜题时，正确性证明会随所解的谜题而不同；如果我们还希望对计算的理解体现在程序结构中，这种情况就显得有些不公允！

\section[论一种程序模型]{论一种程序模型}
\label{eq:program-model}

在拥有一个程序之前，必须先把它构造出来；拥有之后——如果构造它还有任何意义——便要执行它。本节不着重讨论程序构造和程序执行这两项活动，而尝试把程序看作一个静态对象。我们希望把它看作结构高度有序的对象；主要问题是：我们设想什么样的结构，为什么？希望最终得到的程序结构，既便于构造，也便于执行。当然，在思考时，我无法忽略这些过程，但现在并不想讨论它们。具体说，不讨论设计方法——究竟“由外向内”还是反过来——也不讨论实现方面的后果。同样，为了不使问题过于复杂，只讨论顺序程序。

如果只评价程序本身，我想，我的中心要求是：程序应当按我能够理解的方式写下来，应当按我愿意向别人解释它的方式写下来。但若不作进一步限定，这些不过是人人赞成、却空泛的说法。因此，让我看看能否讲得更具体些。

考虑一个非常简单的计算，其中可以区分出三个相继发生的动作：输入数据、处理数据（即真正的计算），以及输出结果。表示这个程序的一种方式，是写成一长串语句：

\noindent
\textbf{begin}

\noindent
\quad\dots

\noindent
\quad\dots

\noindent
\quad\dots

\noindent
\quad\dots

\noindent
\quad\dots

\noindent
\quad\dots

\noindent
\textbf{end}

下面的形式加入一些标号，以便说明：

\noindent
\tabto{10em}\textbf{begin}
\nopagebreak

\noindent
输入开始：\tabto{10em} \quad\dots
\nopagebreak

\noindent
\tabto{10em} \quad\dots

\noindent
处理开始：\tabto{10em} \quad\dots
\nopagebreak

\noindent
\tabto{10em} \quad\dots

\noindent
输出开始：\tabto{10em} \quad\dots
\nopagebreak

\noindent
\tabto{10em} \quad\dots

\noindent
\tabto{10em}\textbf{end}

\noindent
阅读文本时，它会提示我们接下来发生什么。

更好的写法是：

\noindent
\tabto{10em}\textbf{begin}
\nopagebreak

\noindent
输入：\tabto{10em} \quad\textbf{begin} \tabto{14em}\dots
\nopagebreak

\noindent
\tabto{13em} \quad\dots\ \textbf{end};

\noindent
处理：\tabto{10em} \quad\textbf{begin} \tabto{14em}\dots
\nopagebreak

\noindent
\tabto{13em} \quad\dots\ \textbf{end};

\noindent
输出：\tabto{10em} \quad\textbf{begin} \tabto{14em}\dots
\nopagebreak
\nopagebreak

\noindent
\tabto{13em} \quad\dots\ \textbf{end}
\nopagebreak

\noindent
\tabto{10em}\textbf{end}

\noindent
这时，标号与其说是程序文本中某个位置的标记，不如说是后面区域的名称——区域由成对的“\textbf{begin} \textendash  \textbf{end}”括起来——或者说，是计算分解所得三个动作的名称。然而，即便这样看，这三个“标号”仍然只是注释，是提供给有兴趣的人类读者的说明性附加文字；而我希望把它们视为程序不可分割的一部分。无论进一步检查时这些动作会采取什么形式，我希望程序文本中某处明确体现这个事实：计算已分解为三个动作在时间上的相继发生。做到这一点的一种方法，是在某处写出以下三个（抽象）语句的文本次序：
\medskip

{
	\setlength{\parindent}{8em}
	\hspace{-.4em}``输入；处理；输出''
}
\medskip

\noindent
并约定，这三个动作在时间上的顺序，确实由上述文本顺序控制；而三个动作的进一步细化放在“别处”，也许分别描述，但肯定不施加彼此之间的先后顺序。

如果闭式子程序没有在二十多年前发明出来，现在就该发明它了！换句话说，我们又回到熟悉的地方，熟悉到许多读者甚至会觉得受了骗！我倒不这样觉得，因为只要一种经验证的方法仍令人满意，坚持它就绝不该羞愧。但我们应当清楚，希望从中得到哪些好处；必要时，应调整它；最后，还应建立使用它的纪律。因此，让我重新审视子程序这个概念，因为过去一年中，我对它的认识发生了变化。

我最初接触闭式子程序的概念，是在了解 EDSAC 时\footnote{``The Preparation of Programs for an Electronic Digital Computer; with Special Reference to the EDSAC and the use of a Library of Subroutines'', M. V. Wilkes, D. J. Wheeler and S. Gill, Addison-Wesley Press, 1951.}；当时，子程序概念是标准例程库的基础。那个年代，制造硬件是一场伟大的冒险，许多标准例程是以（稀缺的！）存储空间和计算时间换取硬件电路的手段：指令系统不包含除法指令，于是就有除法子程序。奇怪的是，我不记得自己当时曾把子程序理解为一种把给定机器“重新构造”为更合适机器的手段。我也不记得当时曾把子程序看作由用户构思、构造，用来体现其分析的对象；它们更多是供用户使用的标准例程。最终，我主要把它们看成缩短程序长度的手段。但整个程序仍被设想为在单一、均匀的存储器中运行，在一个没有结构的状态空间中活动；整个计算仍被设想为单个处理器执行的单一顺序过程。后来，我讲过许多程序设计课程，忠实地宣讲这套教义，经常解释调用序列怎样传递返回地址，以及子程序怎样先在自身末尾设置“链接”——也就是返回跳转。现在，我更愿意把主程序看成拥有自己的指令计数器，子程序执行完毕后，它只需继续“计数”；肯定不会把那个“休眠的值”视为传给子程序的参数。（不过，这个旧观点仍进入了许多机器的硬件。我们见过这样的机器：子程序跳转把链接存到子程序的“地址零”，并指示从“地址一”恢复取指令。这种安排让可重入代码和递归子程序的实现有些困难。甚至在这十年中，我们仍见到这样的机器：发生程序中断时，把被中断程序的“程序状态”存到与中断相关的位置，而不是与被中断程序相关的位置！）

十年后，ALGOL 60 出现，情形发生变化。我们不再说闭式子程序，而改称“过程”。程序员仍把它们看成缩短程序文本的便利手段；越来越多程序员也开始用它们组织结构，使程序在适应预见到的问题规格变化时，只需替换一个或几个过程体，或在过程调用中改变某些实参。但最主要的新东西，是局部变量的概念。

这体现在两个重要方面。第一是“作用域”概念，即并非所有变量在程序各处都同样可访问：过程的局部变量，在过程体之外无法访问，因为在那里它们无关紧要。过程为了完成自己的任务需要哪些局部变量，是它自己的事；调用它的主程序不必关心。主程序可以——而且必须！——独立于这些局部变量构思，这个事实得到了明智的体现。我们也许对 ALGOL 60 的具体作用域规则有所保留，但应认识到，这是朝正确方向迈出的一大步。

第二个新方面，是过程可以递归使用；更准确地说，允许过程直接或间接调用自己。这种机制的好处曾引起许多激烈争论；据我所见，争论已经平息。反对递归过程的论据始终是效率：不可重入代码可以执行得高效得多。但多道程序设计出现后，灵活分配存储空间又有了新的必要性。如果仍有某些机器执行不可重入代码高效得多，也就是说，使用递归例程的代价过于高昂，那么我冒昧地认为，现在应把这类机器的结构称为有些过时。不过，递归过程迫使我们承认，它的（静态）文本与（动态）激活——也有人称之为“化身”——之间有区别。过程文本是一回事；它这一次所操作的那组局部变量，则完全是另一回事。

到这里，一切都好，但它也有一些缺点；至于你把它们叫作语言上的还是概念上的缺点，我并不在意。进入过程时，“创建”局部变量；退出时，“销毁”它们。正是这种对某个过程化身所含变量生存期的自动控制，让我们可以用栈——后进先出的存储安排——实现（递归）过程。一个化身的局部变量只在该化身期间存在，因此过程无法在幕后把信息从一个化身传给下一个。为克服这一点，引入了“own”概念，但它并未解决问题。第一，递归情况下，own 变量究竟有什么用，很不清楚；第二，无法编写一组共享若干 own 变量的过程。（可以在包围过程声明的外层程序块中声明这些变量来模拟，但作用域规则又让它们可访问的范围太广，不能再算“在幕后”。）我们的结论——这绝不是新结论，也绝非我一人的结论！——是：ALGOL 60 引入的“own”概念考虑不周；我们必须寻找新方式，控制并描述局部变量的生存期、可访问性和身份。

我对过程概念还有另一个不满：它仍主要被当作缩短程序文本的手段（虽然在递归情况下，文本长度可能未知）。过程调用的语义，用著名的“复制规则”描述：过程调用应当理解为一种简写；在语义上，应以过程体文本的一个副本替换它，适当调整标识符、代换参数，然后由执行主程序的同一台机器，执行修改后的文本。它仍是由单一顺序机器执行的单一程序文本的一种表示。而恰恰是这幅单一机器的图景，已不再令我满意。

我希望把主程序看作由它自己专用的机器执行：这台机器具有合适的指令集合，操作合适的变量，并用自身的指令计数器控制执行顺序。如果真有这样一台机器，主程序就能解决我的问题。我希望这样看，是因为它强调：无论这台机器——很可能还只是一台虚拟机器——是否存在，都能讨论并确认主程序的正确性。我不必拥有它，只需知道与正确执行当前主程序有关的那些规格。

对我而言，构想这台虚拟机器，体现了自己的抽象能力。这很像我理解用所谓高级语言写成的程序的方式：不必知道各种操作——例如乘法和下标操作——如何实现，也不必知道那些无关细节，例如最终负责执行程序的硬件采用什么数制。

当然，实际中会发现，这台理想机器并不存在；于是下一项任务——在结构上与原始任务相似——就是编写程序，模拟这台“上层”机器。编写模拟程序时，必须决定数据结构，以提供上层机器的状态空间；还要构造一组算法，分别实现上层机器指令系统中假定存在的各条指令。最后，“下层”机器还可以有一组为自身方便而引入的私有变量，完全在上层机器的领域和作用域之外。但这组程序所针对的机器，很可能仍然不存在；于是下一项工作，就是用再下一层机器的程序模拟它，如此继续，直到最终得到硬件能够执行的程序。

如果成功按以上方式构造程序，就把程序安排成了若干层。每一层都应能够独立理解，只需假定有合适的机器执行它；而每层的作用，是模拟紧邻其上层所假定可用的机器。

为什么采用这个模型？希望从中得到哪些好处？让我列举一下。

\begin{enumerate}[wide, nosep, label=(\arabic*)]
	\item “\nameref{sec:first-example-step-wise-program-composition}”一节记录的经验强烈表明：设置对应于不同抽象层次的各层，是一种适合程序构造的手段。

	\item 我们可以合理期待，许多程序修改现在可以表述为：用一台兼容的机器，替换某一台（虚拟）机器。

	\item 可以期待，这个模型使我们更好地把握程序运行中需要修改时出现的问题。如果给定层次的机器在两条指令之间停止，所有下层机器都处于完全被动状态，可以替换；所有上层机器则必须看作正在执行某条指令，其状态正处于转换之中。在顺序机器中，只有在指令执行之间，才能解释状态；而这幅机器层次图景——每台机器都有自己的指令计数器，“计数它的指令”——似乎更有利于我们在任一时刻决定哪些解释有效。通常的程序设计语言，用一种均匀尺度衡量计算进展，例如以一条语句为“粒度”；在这样的语言中，面对哪些解释在何时有效的问题，我感到有些无从下手。
	
	\item 可以期待，检测到故障之后，这个模型甚至能帮助处理完全或部分恢复问题。（最近，我参与了一个多道程序系统的设计与构造；最令人恼火的事情之一是：当某个存储单元发出奇偶校验警报时，我们完全无法用机械方式估计灾难的范围。唯一能够实现的安全反应，就是立即停止机器；这实在不是一个值得自豪的方案！）
	
	\item 分层机器体系的图景，可以抵消无情运用“分而治之”原则带来的一项危险：不同部件被如此独立地编写，以至于出现重复工作，甚至更糟的情况。一层包含由某台概念上的机器执行的“一组程序”，这一事实强调：这组程序应当共享同一套基本操作。分离任务是好事，但另一方面，还必须把分开的各端重新接起来！
\end{enumerate}

\section[逐步构造程序的第二个例子]{逐步构造程序的第二个例子}

既然程序结构作为分层机器体系的图景已经浮现，我便手痒，想试用它，也就是再写一个程序。这里采用的记法技术不应视为经过深思熟虑的建议；它们是按我的个人喜好选择的，应看作试验的一部分。

问题如下：给定一台行式打印机，由两条命令控制。“NLCR”（New Line Carriage Return，换行并回车）把下一行最左侧的位置定义为“当前可打印位置”；“PRSYM$(n)$”则在当前可打印位置打印由整数参数 $n$ 的值指定的字符，并把紧邻右侧的位置定义为新的当前可打印位置。（讨论中，可以允许行长无限。）我们只用 $n$ 的两个特定值，分别叫作“space”和“mark”。“PRSYM$($space$)$”让当前可打印位置保持空白；“PRSYM$($mark$)$”则打印一个给定的可见字符，比如某种星号。

另外给定两个整数自变量的整数函数，满足
$$
\text{\textbf{for} } 0 \leqslant i < 1000 : 0 \leqslant f_x(i) < 100 \text{ \textbf{and} } 0 \leqslant f_y(i) < 50.
$$

现在，要求编写一个打印 50 行的程序。各行按 $y$ 坐标从上到下编号，范围由 49 到 0；每行的位置按 $x$ 坐标从左到右编号，范围由 0 到 99。在以下关系给出的一千个位置上（如果有重合，实际位置更少），
$$
x = f_x(i) \text{ \textbf{and} } y = f_y(i) \text{ 对某个 } i \text{ 满足 } 0 \leqslant i < 1000
$$

\noindent
必须打印一个标记；纸面上所有其他位置都保持空白。换句话说，以离散参数形式给出一条曲线，我们希望用行式打印机作数字绘图机。

我在口试中大量使用过这个问题。多数学生很快发现：由于没有 OLCR（Old Line Carriage Return，回到原行并回车），也没有“退格”，可打印位置必须按打印命令规定的顺序处理；其次，这个顺序与各个标记的顺序无关——例如把标记按 $i$ 递增的次序编号时。因此，他们很快得出结论：似乎应当利用存储。先遍历一千个 $i$ 值，以方便的方式存储页面图像；随后，再由存储的图像控制页面打印。（更精确地说：我们假定计算机有足够的存储空间，而且计算函数值“$f_x(i)$”和“$f_y(i)$”可能很费时间，所以希望每个 $i$ 值只算一次。）

现在，把这个设计决定记录下来。我建议用以下文本：

COMPFIRST

\textbf{begin}

\quad draw: $\{$build; print$\}$;

\quad\textbf{var} image;

\quad\textbf{instr} build$($image$)$, print$($image$)$

\textbf{end}

这段文档视为最终程序不可分割的一部分，解释如下。它指的是一台叫作“COMPFIRST”的机器（机器名用大写字母，并尽量表达为该机器编写的程序体现了哪类决定）。

下一行给出一个有名字的算法，名为“draw”（假定这是待编写的整个程序的名字，因为程序要“画”一条曲线）。算法表达了两个动作所需的时间顺序：在存储器中建立图像，然后由所存值控制纸面打印。

最后两行给出声明（或声明首部），为运行上述算法的机器各部件命名。第一行说明：名称“image”用于保存页面图像的数据结构；变量“image”是这台机器状态空间的唯一组成部分。指令系统包含两条指令，分别名为“build”和“print”。

继续之前，应指出我们用了一些简写；其中有些究竟明智还是愚蠢，我还不知道。这些简写都与一个事实有关：变量“image”是这种类型的唯一变量。

如果状态空间中有两个图像，我会写成

\noindent
``\textbf{type} image;

\noindent
image \textbf{var} imagel, image2''

\noindent
表示状态空间包含两个变量，分别叫作“imagel”和“image2”；它们有同一个可能值集合，该集合由名为“image”的\textbf{类型}刻画。后续进一步说明\textbf{类型} image 的细节时，说明会同时适用于两个变量。由于这种类型的变量集合只有一个元素，我冒险没有区分这个集合（名为“image”）及其唯一元素（也叫“image”）。COMPFIRST 中的描述，例如“build(image)”，提到“image”时，指变量；以后给出结构细节时，指类型 image。

最后一行包含指令声明，类似于过程首部。一般说，声明包含参数的类型，调用则把变量作为实参。同样，如果参数已由类型唯一确定，这样做似乎就没有必要。为此，我们在声明中提到实参，却在描述算法“draw”的代码中省略“image”。这样，就能把明确写出实参的做法，留给真正使用参数组合自由的情况。

继续之前，我要强调：名为“draw”的这个小算法，可以而且应当看作针对某台机器编写的程序。我们应为这台机器编写手册，说明
\begin{enumerate}[wide, nosep, label=(\arabic*)]
	\item 操作“build”给变量“image”赋值，该值根据函数 $f_x$ 和 $f_y$ 指定将要打印到纸上的图像。

	\item 操作“print”把变量“image”当前值指定的图像打印到纸上。
\end{enumerate}

要看出它确实可以视为某台机器的算法，最容易的方法也许是考虑“draw”的其他算法方案。例如

draw: $\{$print; build$\}$

\noindent
是错误的，因为此时动作“print”没有定义；

draw: $\{$build; build; print$\}$

\noindent
是正确的，却无谓地耗费时间，因为第二次“build”给“image”赋的是它已经具有的值；

draw: $\{$build; print; print$\}$

\noindent
则有意义：它会把图像打印两遍。

现在继续程序设计任务。如果有“COMPFIRST”机器可用，由它执行这个名为“draw”的小程序，就能完成工作。为继续讨论，也为了符合现实，假定我们没有这样一台量身定做的机器。因此，下一项任务——与前一项相似！——就是构造这样一台机器。

目前假定了三个已命名实体：“build”“print”“image”；前两个引用后一个。因此，进一步说明后者的细节，会影响前两者；而且，在没有对“image”的结构作进一步承诺时，很难继续细化“print”。但“build”可以独立进一步细化。因此，我们首先选择“build”作为细化对象。

要说明变量“image”怎样得到一个与一千个标记的正确位置对应的值。作为整体操作，它给一个原先值尚未定义的变量赋值。预见到标记将“每次一个”地加入，我们发现，加入下一个标记，是对“image”已经定义的值进行操作。因此，把全部标记的设置过程看作对某个已定义值——即空白页面对应的值——进行操作，似乎很合适。这个决定得到

CLEARFIRST

\textbf{begin}

\quad build: $\{$clear; setmarks$\}$;

\quad \textbf{instr} clear$($image$)$, setmarks$($image$)$

\textbf{end}

\noindent
其中，动作“clear”给 image 赋予与五十行空白组成的图像对应的值；动作“setmarks”把 image 的初始值调整为已加入曲线的一千个——或更少——标记的值。

CLEARFIRST 同样是一台机器，也可以为它编写其他程序。例如

build: $\{$clear$\}$

\noindent
有意义，但会输出五十行空白；

build: $\{$setmarks; clear$\}$

\noindent
会包含一个未定义的操作；

build: $\{$clear; clear; setmarks$\}$

\noindent
会包含一个多余操作，下面也是如此：

build: $\{$clear; setmarks; setmarks$\}$

\noindent
因为第二次“setmarks”只会添加图像中已经存在的标记，因而不会改变“image”的值。

（关于记法的说明：用“clear”和“setmarks”解释“build”的算法，并没有明确提到“image”。这是因为，除非这台机器确实用到了实参的组合自由，否则我们不想在算法中使用实参记法。

此外，“build”是单参数操作，因此没有为它的形参另引入标识符。我的这个简写，也可能非常不明智。）

计算设计的下一步，是说明怎样依次处理曲线的一千个标记；这一步不需要作任何其他承诺。目前，我建议这样写：

ISCANNER

\textbf{begin integer} $i$;

\quad setmarks: $\{i \coloneq 0$; \textbf{while} $i < 1000$ \textbf{do} $\{$add mark; $i$ \textbf{plus} $1\}\}$;

\quad \textbf{instr} add mark$(i$, image$)$

\textbf{end}

\noindent
这个算法应在一台具有“add mark($i$, image)”指令的机器上理解。该指令根据加入第 $i$ 个标记这件事，改变“image”的值。算法描述了标记的处理顺序，并表明每个标记都恰好处理一次。

但这还不是全部：这里引入了一个新变量“$i$”；算法使用一组与它有关的动作——“$i \coloneq 0$”、\mbox{“$i < 1000$”}和“$i$ \textbf{plus} 1”。如果完全保持一致，我似乎应像“add mark”一样，把它们列在底部，表明它们以后可能还需要进一步说明。但我没有这样做，而把它们按 \textbf{while-do} 子句那样处理。从语言语义角度看，这样单独对待一个隐含理解的类型 \textbf{integer}，似乎不合适；也很难说明，为什么对类型 \textbf{integer} 与类型“image”区别对待，因为在这台机器中，两者都是隐含理解的。

然而，我确实这样做了。我一直为并不存在的机器设计程序，并补上一句：“如果现在有一台机器，包含这里假定的基本操作，工作就完成了。”从逻辑上讲，这是正确的；从实践上看，却是个玩笑，因为我们很清楚，不能假定会有一台通用机器，其指令系统如此符合我们的需要。我们不应忽略——也不应假装忽略！——自己在后续设计阶段提供这些基本操作的责任。现在，我以这种例外方式使用已充分理解的“integer”类型，以及对这类变量定义的操作，是想表明：虽然这些功能必须以某种形式提供，但提供它们不在程序员的责任范围内；程序员也会接受任何合理的实现。

现在，又剩下一个无需承诺其他基本操作、就能进一步细化的基本操作。我们要说明：怎样把处理编号为 $i$ 的标记，表达为处理页面上的某个位置。于是构造一台计算这个位置的机器。

COMPPOS
\nopagebreak

\textbf{begin integer} $x$, $y$;

\nopagebreak

\quad add mark: $\{x \coloneq f_x(i)$; $y \coloneq f_y(i)$; mark pos$\}$;

\quad\textbf{instr} mark pos$(x$, $y$, image$)$

\textbf{end}

\noindent
其中，“mark pos”按照在待打印图像的“$x$”“$y$”坐标处加入标记这一动作，改变变量“image”的当前值。

（注：最后这次细化明确假定，函数 $f_x(i)$ 和 $f_y(i)$ 可以按任意自变量值顺序求值。如果这两千个函数值必须从输入流中，按规定的 $i$ 值顺序成对读入，那么最后两台机器就必须合并为一台。）

到这里，如果不确定仍然相当模糊的“image”类型的结构，我已看不到继续细化的可能。这个值要怎样存储？必须安排 image 类型变量的结构；完全等价地说，必须为其可能值选择一种表示。

在各处讲课时，我曾向不同听众讲述这个程序的几个版本。也许值得指出的是：至少有两次，当我讲到这个阶段，部分听众感到深深不安。例如，他们觉得，我不能声称目前已写出的程序是正确的；他们反对我说

draw: \{build; print; print\}

\noindent
会把同一图像打印两遍。因为在\textit{构造}基本操作“print”之前，我怎么知道它不会通过某种副作用，改变“image”的值？当然，答案是：“print”必须做已声明的事，不应做未声明的事。但随之而来又有别的反对意见：我没有证明表示是唯一的；也许“print”只是一个偏函数，对“image”的某些可能值未定义，等等。对此，答案似乎是：这些担忧本身合理，但应当在适当时候处理，也就是不能在决定表示方式之前就处理。我们的做法显然有一个优势：这么多程序内容，都能独立于“image”类型值的具体表示写下来。迄今所做的，确实似乎是对抽象能力的一种明智利用；这里，抽象掉的是为数据结构“image”选择的具体表示。

但即使现在认定，已到决定“image”类型数据结构的时候，我们仍不必完全定死。面对怎样安排变量结构的问题，可以逐步作决定，正如迄今的算法细化一样。

回想问题的起因：打印基本操作“PRSYM”和“NLCR”迫使计算从上到下，逐行生成图像。让我们试着正视这一事实，把图像看作由行组成的数组。于是得到下一层：

LINER

\textbf{begin integer} $j$;

\quad image: $\{$\textbf{array} line$[0: 49]\}$;

\quad print: $\{j \coloneq 49$; \textbf{while} $j \geqslant 0$ \textbf{do} $\{$lineprint$($line$[j])$; $j$ \textbf{minus} $1\}\}$;

\quad clear: $\{j \coloneq 49$; \textbf{while} $j \geqslant 0$ \textbf{do} $\{$lineclear$($line$[j])$; $j$ \textbf{minus} $1\}\}$;

\quad mark pos: $\{$linemark$($line$[y])\}$;

\quad \textbf{type} line;

\quad \textbf{instr} lineprint$($line$)$, lineclear$($line$)$, linemark$(x$, line$)$

\textbf{end}

倒数第二行中，引入了名为“line”的类型。再次说明，类型被看作一个可区分的值集合；这种类型的变量，在任一时刻都可以取集合中的一个值。代码第一行表示：“image”类型由 50 个“line”类型元素构成的数组组成，编号从 0 到 49。由于它是唯一一个由这种类型组成的类型，我们再次没有引入新标识符；这样做是否明智，尚且不论。

然后，把原来理解为对“image”操作的“print”“clear”“mark pos”，翻译为用行操作表达的算法。在这些算法代码中，真正的实参指定是哪一行；描述最后给出指令列表，表明动作操作的是“一行”：给了类型，却没给参数。

这一层带来了若干新特点。首先，如同解释“image”时所做的，我们把数据类型的结构细化，与算法细化——例如对“print”“clear”“mark pos”的细化——置于非常相近的地位。在此之前，我们的方法还可以看作试图为“子程序化”建立纪律——请读者原谅这个难听的说法！——现在却看到，这样概括只涵盖了我们努力的一半，因为我们也试图把相似技术用于数据结构。第二，之前的机器只解释一个实体——指令或数据类型；“LINER”却解释了一整组实体。关键在于，我们尝试为每层对应一个独立设计决定。这里的决定是：从此以后，以行来理解图像；因此，所有对图像本身的操作，都必须翻译成对其各行的操作。图像已经通过行得到具体解释，在下一层不再出现；唯一还需要处理的不寻常类型是“line”，接下来就处理它。请注意：下一层处理的是行；行是用来组成图像的，这件事已不再相关！

表示一行，有许多不同可能。例如，列出应打印标记位置的 $x$ 坐标（可能按 $x$ 值递增排序）；使用 100 个元素的布尔数组，每个元素指示该行相应位置是否应有标记；或者使用 100 个元素的整数数组，每个元素存放相应可打印位置所需的 PRSYM 参数值，即“mark”或“space”。最后一种表示便于扩展到在同一幅图中打印不同曲线、采用不同标记的情况，因此选它。

由此得到

\noindent
LONGREP

\noindent
\textbf{begin integer} $k$;

\noindent
\quad line: $\{$\textbf{integer array} sym$[0 : 99]\}$;

\noindent
\quad lineprint: $\{k \coloneq 0$; \textbf{while} $k < 100$ \textbf{do} $\{$PRSYM$($sym$[k]$); $k$ \textbf{plus} $1\}$; NLCR$\};$

\noindent
\quad lineclear: $\{k \coloneq 0$; \textbf{while} $k < 100$ \textbf{do} $\{$sym$[k]\ \coloneq$ space; $k$ \textbf{plus} $1\}\};$

\noindent
\quad linemark: $\{$sym$[x]$ $\coloneq$ mark$\}$

\noindent
\textbf{end}

不过，这种实现会把一行中最右侧标记右边的部分全用空格填满：这就像写信时，为了另起一段，不停敲空格键，直到打字机的铃响！

下一版本省去了多余的 PRSYM 命令，甚至让“line”类型变量中那些无需定义的元素保持未定义。每行关联一个计数器“$f$”，给出该行需要多少条 PRSYM 命令。清空一行，现在简化为把“$f$”置零！

\noindent
SHORTREP

\noindent
\textbf{begin integer} $k$;

\noindent
\quad line: $\{$\textbf{integer} $f$; \textbf{integer array} sym$[0 : 99]\}$;

\noindent
\quad lineprint: $\{k \coloneq$ 0; \textbf{while} $k < f$ \textbf{do} $\{$PRSYM$($sym$[k])$; $k$ \textbf{plus} $1\}$; NLCR$\}$;

\noindent
\quad lineclear: $\{f \coloneq 0\}$;

\noindent
\quad linemark: $\{$sym$[x] \coloneq$ mark;

\noindent
\quad \textbf{if} $f \leqslant x$ \textbf{do} $\{k \coloneq f$; \textbf{while} $k < x$ \textbf{do} $\{$sym$[k] \coloneq$ space; $k$ \textbf{plus} $1\}$;

\noindent
\quad \quad $f \coloneq x + 1\}\}$

\noindent
\textbf{end}

\noindent
\textit{后来补充的说明。}

上面的程序，基本就是我曾向至少五批听众展示的程序。现在，两个月之后，在从容思考正确性证明时，我忽然意识到：上面的“linemark”算法暴露了我过去的习惯。与以下方案相比，它是一段糟糕的代码：

linemark: $\{$\textbf{while} $f \leqslant x$ \textbf{do} $\{$sym$[f] \coloneq$ space; $f$ \textbf{plus} $1\};$
	
\quad sym$[x] \coloneq$ mark$\}$

\noindent
这个版本保证：每当执行“sym$[x] \coloneq$ mark”时，总满足关系“$x < f$”；第一行的职责，正是确保这一点。请读者尝试理解两个 linemark 版本，并比较各自的推理方式。然后，你就会赞同我的判断：原来的版本很糟。

第二个版本之所以突然浮现在脑中，是由于以下观察。条件子句
$$
\text{``\textbf{if} } B \text{ \textbf{do} } S\text{''}
$$

\noindent
在程序中有两种用法。一种是执行语句 $S$ 后，$B$ 仍然为真；另一种是可以保证，执行语句 $S$ 会使 $B$ 不再为真。后一种情况下，条件语句的职责，是保证执行之后 $B$ 不成立。此时，它实质上是
$$
\text{``\textbf{while} } B \text{ \textbf{do} } S\text{'',}
$$

\noindent
的简写；后者具有使 $B$ 不再为真的性质（前提是它终止）。但要说明这一简写合理，还必须另行证明：被重复的语句至多执行一次。（在“\nameref{sec:first-example-step-wise-program-composition}”中，我们没有费心在第 2b4(4) 层引入这个简写，而是写成
$$
\text{``\textbf{while} “ord 太小” \textbf{do} “将 ord 加一”'';}
$$

\noindent
那里用一个条件子句就足以完成工作！）

\section[论我们已经取得的成果]{论我们已经取得的成果}

思考所设想的程序结构时，我常用一个比喻：把程序看作由一颗颗珍珠串成的项链。

我们用若干层次描述程序，每层都含有对较高层次假定可用实体的“细化”。细化可以是动态的——算法——也可以是静态的——数据结构——由适当的机器来理解。我用“珍珠”一词，指这样一台机器，连同它包含的细化。

上一个程序是一条由六颗珍珠串成的项链，顺序或者是

COMPFIRST

CLEARFIRST

ISCANNER

COMPPOS

LINER

LONG REP

\noindent
或者是

COMP FIRST

CLEARFIRST

ISCANNER

COMPPOS

LINER

SHORTREP。

LONGREP 和 SHORTREP 是两颗不同珍珠：它们把“上端面”的同一些概念，解释为“下端面”的同一些概念；不同的只是具体细化。因此，它们相当于为同一机器、完成同一工作而写的不同程序方案。

修改程序，将被视为把原项链的一颗或几颗珍珠，替换为另一颗或几颗珍珠。珍珠是构成程序的独立单位。构造一个程序——作为相关程序类中的一员——现在被看作两个阶段：先制作珍珠，数目多于严格所需；再从中选择，串成一条合适的项链。

这种两阶段方法有许多理由。设计程序时，必须考虑许多、许多不同的程序方案；程序完成后，还必须把它改成某个其他方案。只要仍把程序视为程序设计语言基本符号的线性串，并相应地把程序修改看作这一层次的文本操作，那么每项修改都必须放在“该语言能写出的所有程序——正确或错误的程序！”这个全集中理解。如此一来，程序修改是一项风险极高的操作，也就不足为奇了！基本符号作为描述这种修改的单位，太小，而且没有意义。珍珠则体现一个独立设计决定，或者原始问题说明中一个独立的方面；我们希望它成为这类修改的自然单位。

换一种方式说明同一论点：随着 ALGOL 60 诞生，人们发现，语法是表达程序文本结构的强大手段。（语法受到如此推崇，以至于许多领域工作者把计算机科学等同于语法分析！）不过，一个事实多少被忽略了：通过语法表达结构，结构只能非常间接地给出；必须在线性的基本符号序列上应用语法分析算法，才能导出它。意识到许多程序修改并不影响结构的大部分之后，这就显得令人痛苦：费力地重新分析修改后的文本，最后又得到原来的结构！我强烈觉得，上下文无关方法表示结构的适用程度，被严重高估了。（我身边有人向我展示过以下 ALGOL 60 程序错误。一个程序在具备完整检查功能的实现中产生错误输出。除了程序文本，这种实现还要求在最后一个“\textbf{end}”之后，加上最终的“\textbf{progend}”；其他任何位置都不允许这个附加符号，从而能确认“\textbf{begin} \textendash  \textbf{end}”的配对正确。结果发现：

\begin{enumerate}[wide, nosep, label=(\arabic*)]
	\item 程序某处漏了一个字符串结束引号；
	
	\item 后面某处又漏了一个字符串开始引号；
	
	\item 由此得到的程序，其“\textbf{begin} \textendash  \textbf{end}”结构在语法上正确；
	
	\item 两个遗漏位置之间声明的标识符，只在这两个位置之间使用，因此连上下文相关的检查也无法报警。
\end{enumerate}

\noindent
我本来就怀疑上下文无关方法是否适合表达宏观结构，因此看到这个错误时，倒感到十分欣喜！）

越思考珍珠，我越觉得：如果承认自己有责任为一个大型程序考虑，比如说，一千个可能版本，那么只有类似珍珠的东西，才是出路。不能指望程序员把这一千个版本从头写起，彼此独立地完成。我看到的唯一办法，是用组合技术产生这种潜在的多样性：制作多于单条项链所需的珍珠，例如制作 250 颗，单条项链需要 200 颗；再选择某组珍珠串成项链。我看不到其他可行方式。另一种通过组合产生大量变化的机制，是排列；但我们用不了它，因为最终项链必须彼此吻合，而给定珍珠后，为得到吻合项链应采用的串接顺序，已相当明确。此外，即使顺序并非唯一，允许改变的那些顺序也相当无关紧要！

珍珠也为一个“不完整”的程序——仅含项链的上半段——赋予明确地位：它可以看作由某台适当机器执行的完整程序；项链的下半段则给出该机器的可行实现。因此，上半段的正确性可以独立于下半段的选择来确认。在相邻两颗珍珠之间，可以作一道“切口”；切口处是一份机器手册，该机器由下半段提供，供上半段所代表的程序使用。这份手册成为上下两部分之间的接口。相比把数据表示看作操作之间的接口，我们觉得这种接口更有帮助；特别是在程序需要调整时，更有助于确保所需的组合自由。

另一点意见涉及概念沿项链的有效范围。例如，“image”概念在最上面的珍珠“COMPFIRST”中引入，在倒数第二颗珍珠“LINER”中得到具体解释，此后不再出现。如果现在认定，设想中的程序对存储空间要求太高，无法引入变量“image”，就面临一次重大修改：必须用其他珍珠替换最上面的五颗，因为这就是“image”概念的有效范围！最底下那颗珍珠——“LONGREP”或“SHORTREP”——却可以保留。（我举这个例子，是为了说明，珍珠替换绝不限于替换最底下一颗。）

关于概念沿项链的有效范围，我想请大家留意一个观察。当我第一次发现它时，非常激动。（回头看，它相当显然，而这恰恰说明值得明确讲出来。）每颗珍珠对应“一个独立设计决定”，因此，珍珠沿项链的次序也就意味着设计决定的次序。能改变这个次序吗？能，不过那样就会得到不同的珍珠。作为试验，我遵照了一个众所周知的建议：面对两个基本操作——这里是“build”和“print”——立即确定它们的接口——这里是“image”——这样，两个基本操作就能互不依赖地细化。我照做了，得到以下项链形式：

COMPFIRST

LINER'

CLEARFIRST'

ISCANNER'

COMPPOS'

SHORTREP

\noindent
（中间四颗珍珠加上撇号，表示它们与原来的珍珠不同，虽然体现的设计决定相同。）得到的程序却乱得多。为什么？

沿着项链，可以标出每个概念的有效范围；这些范围当然彼此重叠。我们可以把它们看成一根根线，整个解释由这些线绞合而成，就像一条“逻辑绳索”。混乱的版本用更多、而且有些更长的线绞出逻辑绳索：绳索更粗，整个构造的相互联系更紧密。这一观察令我激动，因为它非常有说服力地表明——至少对我如此！——优雅、清晰等品质，确实有显著的数量方面的含义。（莫扎特明白这一点：他的许多令人屏息的作品，简单得具有欺骗性，看起来几乎是用什么都没有的材料写成的！）

也可以用更技术性的说法表达这一观察。带撇号的版本很乱，因为过早用行来具体解释图像，使图像概念不再出现在随后层次中，迫使我们也用行来解释“CLEARFIRST”“ISCANNER”和“COMPPOS”；实际上，它们仍可以用图像解释，也就是独立于将为图像选择的表示方式。换句话说，原版本比带撇号的版本，更明智地利用了抽象能力。不依赖特定表示的珍珠越多，程序就越容易适应变化，也越容易理解，因为这一组珍珠可以在较高的抽象层次上理解。经验似乎表明，可适应性和清晰性的目标，已经有了一些具体内容；而且，两者本质上相辅相成。这很令人鼓舞，虽然并不令人意外。

这也让我——至少让我——对当前努力的范围看得更清楚。无论最终发展出什么，它都不会是一个 General Problem Solver：既不是机械的问题求解器，也不是专为人类问题求解者写的一套通用方案。但当一个人成功找到解法时，它也许能帮助他认识解法在“优雅”方面的各种品质；因此，也许能给他一些指引。

\section[论分组与执行顺序]{论分组与执行顺序}
\label{sec:grouping-and-sequencing}

虽然我们正考虑一种明确承认抽象层次体系的程序设计工具，但本节同样适用于以今天所理解的程序设计语言编程，也就是在恒定语义层次上编程。（而且，本节的启示很可能不限于程序设计这个狭窄领域，因为它似乎涉及一般的问题求解。）

我将用两个例子说明，同样都曾用于口试。第一个例子得自 Niklaus Wirth。

问题是：编写程序，生成由 $0$、$1$、$2$ 组成的非空序列；序列中不得有两个相邻、非空且逐元素相等的子序列。按照字典序生成这些序列，直到生成一个长度为 100——即 100 个数字——的序列。程序员可以利用一个已知事实：确实存在满足条件、长度为 100 的序列。待生成序列列表的开头是：

0

01

010

0102

01020

010201

0102010

0102012

\dots

每个解，除第一个以外，都是在较早的某个解上扩展一个数字得到的；因此，这自然是一个直接的回溯算法。

我们寻找“合格”序列，并假定有一个基本操作，可以检查试探序列是否合格。如果合格，就打印它，并在末尾加一个零，得到下一个试探序列；如果不合格，就对它执行“递增”操作，得到下一个试探序列：移除所有值 $= 2$ 的末尾数字，然后把最后剩下的数字加 1。（“在末尾加零”和“递增”保证试探序列按字典序生成；解是其中的一个子集，因此也会按字典序打印。）算法首先检查以下试探序列；带星号的会因“不合格”而被排除：

\noindent
\tabto{\parindent}0

\noindent
*\tabto{\parindent}00

\noindent
\tabto{\parindent}01

\noindent
\tabto{\parindent}010

\noindent
*\tabto{\parindent}0100

\noindent
*\tabto{\parindent}0101

\noindent
\tabto{\parindent}0102

\noindent
\tabto{\parindent}01020

\noindent
*\tabto{\parindent}010200

\noindent
\tabto{\parindent}010201

\noindent
\tabto{\parindent}0102010

\noindent
*\tabto{\parindent}01020100

\noindent
*\tabto{\parindent}01020101

\noindent
*\tabto{\parindent}01020102

\noindent
*\tabto{\parindent}0102011

\noindent
\tabto{\parindent}0102012

\dots

我发现，多数学生倾向于写成以下结构的程序：

\hspace{-.4em}``将试探序列设为单个零；

\textbf{repeat if} 合格 \textbf{then}

\quad \textbf{begin} 打印试探序列；在试探序列末尾加零 \textbf{end}

\quad\quad\quad \textbf{else}

\quad 递增试探序列

\textbf{until} 长度 $= 101$''

虽然这样的程序产生正确输出，我们仍然可以——按我的看法，也应该——提出反对意见。第一点涉及停止条件：当长度为 100 的解已经打印时，\textit{我们}了解算法，能够推断，此后试探序列会首次达到长度 $= 101$，于是把它当作停止条件。但确定停止条件的这种方式相当间接、曲折。它有多曲折，那些没有意识到生成了一个不必要试探序列的学生就清楚地展示出来了：他们给试探序列声明了 100 个元素的数组，而不是 101 个。第二点是，“递增试探序列”从来不会增加长度；所以，排除某个试探序列之后，还会作一次多余的长度检查。（使用这个例子考试时，我在课上还没有非常明确地强调问题求解原则，因此失望并不太深。从某种意义上说，我很高兴观察过这些考试，因为它促使我尽可能寻找、表述并教授问题求解原则。）

下面这个程序没有上述问题。它把空序列视为一个不需打印的虚拟解。在同样的细节层次上，结构如下：

\quad \hspace{-.4em}``将试探序列置为空；

\quad \textbf{repeat} 在试探序列末尾加零；

\quad \quad \textbf{while} 不合格 \textbf{do} 递增试探序列；

\quad \quad 打印试探序列

\quad \textbf{until} 长度 $= 100$''

这里，长度就是已打印解的长度——如果已有解打印——因此避免了为停止条件作曲折推理。也不会生成最后那个多余、永远不作检查的试探序列。由于两个循环彼此嵌套，多余的长度检查也不再发生。以效率为主要标准的人，很可能最能接受最后这个理由。我很重视可理解性，因此更喜欢后一个程序：可以把它解释成以下程序结构的进一步细化。

\quad \hspace{-.4em}``将序列置为空；

\quad \textbf{repeat} 将序列变换为下一个解；

\quad \quad 打印序列

\quad \textbf{until} 长度 $= 100$''

这个更抽象的程序，只关心那些是解的序列。在这一描述层次上，可以忽略：从一个解转到下一个解的过程中，会经过一系列最后证明失败的试探解。

第二个例子得自 Joe Weizenbaum。对于给定的正整数 $n$，编写程序，确定最小数 $s$，使它至少可以用两种非平凡地不同的方式，分解为两个 $n$ 次幂之和。

\noindent
（\textbf{for}\tabto{2em}$n = 1$\tabto{6em}$s = 2 = 0^1 + 2^1 = 1^1 + 1^1$

\tabto{2em}$n = 2$\tabto{6em}$s = 25 = 0^2 + 5^2 = 3^2 + 4^2$

\tabto{2em}$n=3$\tabto{6em}$s=1729=1^3 +12^3 =9^3 +10^3$

\tabto{2em}$n = 4$\tabto{6em}$s = 635318657 = 59^4 + 158^4 = 133^4 + 134^4$）

第一次在口试中使用这个例子时，学生花了二十分钟，才多少熟悉这个问题。然后，他勾勒了一个搜索算法；经过修补后，它确实能找到一个有多种“两项 $n$ 次幂之和”分解方式的数。但他无法证明，算法产出的 $s$ 值就是最小值。（事实上，到那时，他一直忽略了问题说明的这一部分。）

于是，他重新组织思路，写成如下形式的程序：

\quad \hspace{-.4em}``\textbf{integer} $s$, $k$;

\quad $s:= 1$;

\quad \textbf{repeat} $s\coloneq s + 1$;

\quad \quad $k\coloneq$ “$s$ 可以分解为两个 $n$ 次幂之和

\quad \quad\quad 的方式数”

\quad \textbf{until} $k > 1$

\noindent
结果得到一个效率低得毫无希望的算法。他的错误，是过早决定依次考察自然数；这些数中，绝大多数根本不能作这种分解。若想到，所寻找的是满足额外性质的最小可分解数，就能得到一个算法，其初步草图可以是

\quad \hspace{-.4em}``\textbf{integer} $k$, $s$, $t$;

\quad $t \coloneq 1$（以及其他初始化）；

\quad \textbf{repeat} $s \coloneq$ “大于 $t$ 的最小可分解值”；

\quad \quad $k \coloneq$ “达到上述最小值的方式数”

\quad \quad $t \coloneq s$

\quad \textbf{until} $k > 1$''

存储一组三元组——数对及其对应的 $s$ 值——使每次 $s$ 大于 $t$ 的最小数对必在其中，并在每次增加 $t$ 时调整这组数据，就能得到一个效率高出若干数量级的程序。$t$ 从一个可分解值，跳到下一个可分解值。程序设计——或一般的问题求解？——就是明智地推迟决定和承诺！

\section[更详细的设计考虑]{更详细的设计考虑}

前面的几节——尤其是“\nameref{sec:first-example-step-wise-program-composition}”——招来了批评，说我把设计过程过分简化，几乎到了不诚实的地步。我认为，这种批评并非全无道理。为补救这一点，我将更详细地处理两个例子。第一个例子是我自己想出来的，曾在几次口试中尝试，最后在“程序设计艺术导论”课程结束时，用于课堂。我把问题交给五十名学生，由我主持讨论，大家一起在 90 分钟内解决了它。

考虑一个字符集，包括字母、空格（sp）和句点（pnt）。单词由一个或多个、但至多二十个字母组成。输入文本含一个或多个单词，单词间以一个或多个空格分隔，最后以零个或多个空格、随后一个句点结束。通过返回字符值的函数 RNC（Read Next Character，读取下一个字符），读取输入文本，从第一个单词的第一个字母开始，到结束句点为止，两者都包括在内。输出文本则使用基本操作 PNC($x$)（Print Next Character，打印下一个字符）生成，参数取字符值。如果程序只需复制文本，那么以下程序就够了（假定字符类型变量可用）：

\quad \textbf{char} $x$;

\quad \textbf{repeat} $x \coloneq$ RNC; PNC($x$) \textbf{until} $x =$ pnt

但在这个例子中，文本必须作以下变换：

\begin{enumerate}[wide, nosep, label=(\arabic*)]
	\item 输出文本中，相邻单词之间必须只有一个空格。
	\item 输出文本中，最后一个单词之后必须只有一个句点。
	\item 从左到右——即反复调用 RNC 扫描的次序——把单词编号为 $0,\ 1,\ 2,\ 3,\ \dots$。偶数序号的单词直接复制，奇数序号的单词则必须倒序打印字母。
\end{enumerate}

例如，用“\textendash”表示空格，输入文本

\quad ``this\textemdash{}is\textemdash{}a\textemdash{}silly\textemdash{}program\textemdash{}.''

\noindent
必须变成

\quad ``this\textendash si\textendash a\textendash yllis\textendash program.''

我诚挚地邀请读者，在继续阅读之前，自己尝试编写这个程序，并记录思考过程，以便与下文比较。（有经验的程序员应当远用不了 90 分钟！）

非空输入文本的长度未知，这提示我们使用以下结构：

\quad prelude;

\quad \textbf{repeat} something \textbf{until} ready;

\quad coda

\noindent
但立即出现一个问题：“每次执行‘something’时，究竟处理多大一部分？”有人提出四种建议：

\begin{enumerate}[wide, nosep, label=(\arabic*)]
	\item 输入文本中的单个字符。
	\item 输出文本中的单个字符。
	\item 两段文本中的一个单词。
	\item 两段文本中的两个连续单词。
\end{enumerate}

前两个建议很快就被否决，没有多少明确说明的理由，虽然——还是因为？——理由并不难给出。（第一种不理想，因为读入下一个字符后，能产生多少输出，变化很大；第二种也有类似问题。此外，一般而言，循环中嵌套循环的程序更清楚，这提示我们寻找较大的处理单位。）听众否决了第四个建议，理由是结束句点既可能跟在偶数个单词之后，也可能跟在奇数个单词之后。为了明确选择第三个建议，我们在黑板上写下：

\quad prelude;

\quad \textbf{repeat} 处理下一个单词 \textbf{until} 已读到句点；

\quad coda

大家满意的是，这个程序整齐地表达了：输出单词的处理次序，与相应输入单词的读取次序完全一致。但它没有表达，一半的单词要倒序打印。指出这一点后，他们很快为此引入一个状态变量。最初有人建议计数已处理单词的个数，再依计数的奇偶性选择处理方式；我略微犹豫了一下，就足以让他们发现，一个布尔变量便能满足需要。于是决定，“prelude”中应包含
$$
\text{``forward} \coloneq \text{\textbf{true}''}
$$

\noindent
而在“处理下一个单词”中，按“forward”当前值指定的顺序打印之后，应执行
$$
\text{``forward} \coloneq \text{\textbf{non} forward''}
$$

在操心单词分隔的细节之前，他们就引入了“forward”；这让我非常欣慰。接下来，分隔成为他们的关注。要意识到继续细化“处理下一个单词”，必须精确说明每次执行被重复语句时，读哪些输入字符、打印哪些输出字符，他们花了更长时间。事实上，我不得不向他们提出这个问题，然后又问：在两段文本中，哪一段更自然地呈现出这种分组？他们选择输出文本，并作了以下分组，竖线表示分隔：

\quad $\mid$this\textendash $\mid$si\textendash $\mid$a\textendash $\mid$yllis\textendash $\mid$program.$\mid$

\noindent
也就是以一个单词加上适当结束符为单位。于是，我问输入字符应如何对应分组。当他们注意到结束符之后，其中一人指出：只有“看见”下一个输入单词的第一个字母，程序才“知道”当前输出单词后面应接一个空格。他们随即从右到左，提出以下输入分组：

\quad $\mid$this\textemdash{}i$\mid$s\textemdash{}a$\mid$\textemdash{}s$\mid$illy\textemdash{}p$\mid$rogram\textemdash{}.$\mid$，

\noindent
然后，我保持沉默，让他们盯着这些符号分组，直到有人发现：第一个输入单词的字符分组是个例外，不够优雅；应当改为

\quad t$\mid$his\textemdash{}i$\mid$s\textemdash{}a$\mid$\textemdash{}s$\mid$illy\textemdash{}p$\mid$rogram\textemdash{}.$\mid$，

\noindent
也就是说，第一个单词的第一个字符，应在前置部分读取。于是又引入一个变量，得到

\quad \textbf{boolean} forward; \textbf{char} $x$;

\quad forward $\coloneq$ \textbf{true}; $x \coloneq$ RNC;

\quad \textbf{repeat} 处理下一个单词；

\quad \quad forward $\coloneq$ \textbf{non} forward

\quad \textbf{until} $x =$ pnt

\noindent
第二行就是前置部分；这时，也已决定收尾部分可以为空。

上述阶段是在最初 45 分钟内达到的，然后我们休息，喝咖啡。我个人觉得问题已经解决，接下来不过是常规工作。但事实表明，听众还不够熟练，完成程序又花了 45 分钟。

他们一致决定引入一个

\quad \textbf{char array} $c[1:20]$

\noindent
来存储单词的字母。（没人发现，读入字母并倒序打印，可以用递归例程完成！）实质上，需要做四件事：读取单词字母，打印单词字母，读取足够多的内容以决定应打印哪个结束符，以及打印结束符。我没有列出这四个动作，也没有让他们明确决定如何分组或组合这些动作。听众决定先完成全部读取，再完成全部打印。（这对他们来说，几乎不是一个有意识的决定。）

试图细化读取和打印过程时，他们碰到一项未曾预料的障碍：他们迟迟才发现——至少让我意外——仍需定义读取与打印之间的接口，用它传递待处理单词，哪怕这个接口多么显然。过了很久，才有人说清楚：当从左到右读单词时，$c[i]$ 应等于第 $i$ 个字符。也许一半听众觉得没有必要如此小题大做，但他们又花了同样长的时间，才发现单词长度也需要某种表示。没有人建议用存储一个结束符来表示；他们引入独立整数“$l$”，计数单词的字母。他们决定，第一个单词“this”应表示为

\quad $c[1] = $``$t$'', $c[2] = $``$h$'', $c[3] = $``$i$'', $c[4] = $``$s$'' 且 $l = 4$

他们仍难以写出读取循环。直到我反复说：“我们已经决定，‘$l$’将表示数组中所存单词的字母个数。”他们才得到读取过程的开头：

\quad $l \coloneq 0; $

\quad \textbf{repeat} $l \coloneq l + 1;$ $c[l] \coloneq x;$ $x\coloneq$ RNC \textbf{until} $x =$ sp \textbf{or} $x =$ pnt

\noindent
（初稿漏了“\textbf{or} $x =$ pnt”，但很快补上。）写到黑板上后，他们没有多少犹豫，就完成了读取过程：

\quad \textbf{while} $x =$ sp \textbf{do} $x\coloneq$ RNC

转向打印过程时，他们更有成效了。显然，读取过程已向他们展示接口的用途，各方面的建议纷纷出现。我之前从没向他们讲过那个两难选择（参见第 \pageref{pg:dilema} 页）：是用分支子句在两个重复之间选择，还是用重复子句反复执行一个分支语句。我等着它自然出现；它果然出现后，我就向他们指出。接着，一个学生让我意外：他建议借助更多变量，把两个循环映射到同一形式。于是引入三个整数“$k$、inc、term”，字母打印变为

\quad \textbf{if} forward \textbf{then begin} $k \coloneq 0$; inc $\coloneq + 1$; term $\coloneq l$ \textbf{end}

\tabto{5.8em}\quad \textbf{else begin} $k \coloneq l + 1$; inc $\coloneq - 1$; term $\coloneq 1$ \textbf{end};

\quad \textbf{repeat} $k \coloneq k + \text{inc}$; $\text{PNC}(c[k])$ \textbf{until} $k =$ term

\noindent
随后是

\quad \textbf{if} $x =$ pnt \textbf{then} PNC(pnt) \textbf{else} PNC(sp).

这样，就得到了以下程序：

\quad \textbf{boolean} forward; \textbf{char} $x$; \textbf{char array} $c[1:20]$; \textbf{integer} $l$, $k$, inc, term;

\quad forward $\coloneq$ \textbf{true}; $x \coloneq$ RNC;

\quad \textbf{repeat} $l \coloneq 0$;

\quad \quad \textbf{repeat} $l\coloneq l + 1$; $c[l]\coloneq x$; $x \coloneq$ RNC \textbf{until} $x =$ sp \textbf{or} $x =$ pnt;

\quad \quad \textbf{while} $x =$ sp \textbf{do} $x\coloneq$ RNC;

\quad \quad \textbf{if} forward \textbf{then begin} $k\coloneq 0$; inc $\coloneq + 1$; term $\coloneq l$ \textbf{end}

\quad \quad\quad \textbf{else begin} $k\coloneq l + 1$; inc $\coloneq -1$; term $\coloneq 1$ \textbf{end};

\quad \quad \textbf{repeat} $k\coloneq k + \text{inc}$; PNC($c[k]$) \textbf{until} $k =$ term;

\quad \quad \textbf{if} $x = \text{pnt}$ \textbf{then} PNC(pnt) \textbf{else} PNC(sp);

\quad \quad forward $\coloneq$ \textbf{non} forward

\quad \textbf{until} $x =$ pnt

收录这一节，并非因为问题特别有趣。恰恰相反，我倒想说，它也许太简单，未必适合检验一种有条理的程序构造方法。之所以收录，是因为它真实地记录了我亲眼看到的课堂过程。可以把它看作对一个常被问到的问题的部分回答：程序设计风格，我究竟能教到什么程度？（我从未把《结构化程序设计札记》用于教学；它主要写给自己，也许还写给同行。本节的课堂试验发生在“程序设计艺术导论”课程结束时；那门课另有讲义，包含练习等内容。写下这些文字时，选课学生还没参加考试，因此教学究竟成功到什么程度，对我仍是一个未定问题。他们喜欢这门课；我还听说，他们把我的程序叫作“逻辑诗”，因此我抱着很好的希望。）

\section[八皇后问题]{八皇后问题}

这最后一节改编自我的“程序设计艺术导论”讲义。这个例子与许多其他好例子一样，得自 Niklaus Wirth。加入本节有两个理由。

第一，再次努力更充分地呈现发明的过程。（实际上，我从学生尚不熟悉回溯概念的地方开始，目标是在推进过程中发现它。）

第二个理由更重要：这一节讨论递归作为程序设计技术的使用。前面几节——特别是“\nameref{eq:program-model}”——重新审视了子程序概念；在那里，它体现的是我也称作“操作抽象”的东西。在主程序与子程序的关系中，可以很清楚地区分两个语义层次。主程序层次上，子程序代表一个原子动作；使用它，是因为“它为我们做什么”，而“它如何工作”在这个层次上无关紧要。子程序体层次上，关心的是它如何工作，但可以——也应该——抽象掉它怎样被使用。对递归过程的设计者来说，“它做什么”与“它如何工作”这两个语义层次，无法这样清楚分离。因此，设计递归例程需要一种不同的思维能力，这也说明本节值得收入本稿。递归过程必须在同一个语义层次上理解和构思；在这个意义上，它更像一种执行顺序机制，可与重复子句相比。

要求编写程序，生成在 $8\times 8$ 格国际象棋棋盘上放置八个皇后的所有配置，使任何一个皇后都不能吃掉其他皇后。这意味着，在所求配置中，任意两个皇后都不能在同一行、同一列或同一条对角线上。

我们没有生成所有这些配置的操作，恰恰需要构造的就是它。处理这类问题，有一个非常通用的方法（参见“\nameref{sec:grouping-and-sequencing}”），如下。

把待生成的配置集合叫作集合 $A$。寻找一个配置集合 $B$，使它具有以下性质：

\begin{enumerate}[wide, nosep, label=(\arabic*)]
	\item \label{en:first-item-eight-queen}
	集合 $A$ 是集合 $B$ 的子集。
	
	\item \label{en:second-item-eight-queen}
	给定集合 $B$ 的元素，不太难判断它是否也属于集合 $A$。
	
	\item \label{en:third-item-eight-queen}
	能够构造一个操作，生成集合 $B$ 的全部元素。
\end{enumerate}

借助性质 \ref{en:third-item-eight-queen} 中对集合 $B$ 元素的生成操作，便能依次生成 $B$ 的全部元素；然后应用判据 \ref{en:second-item-eight-queen}，决定跳过还是交付它们，从而生成集合 $A$ 的元素。由于性质 \ref{en:first-item-eight-queen}，这个算法会产生集合 $A$ 的\textit{全部}元素。

需要说明三点。

\begin{enumerate}[wide, nosep, label=(\arabic*)]
	\item 要使整个方法有意义，集合 $B$ 就不能与集合 $A$ 相同；而它必须把 $A$ 包含为真子集，所以必定比 $A$ 大。不过，从效率考虑，最好让 $B$“尽可能小”：元素越多，需要按判据 \ref{en:criterion-2-eight-queen} 跳过的元素也越多。
	
	\item \label{en:criterion-2-eight-queen}
	应当寻找一个应用代价小的判据；至少发现 $B$ 的某个元素\textit{不}属于 $A$，平均而言应当代价很小。这也是效率要求，因为可以预期 $B$ 会比 $A$ 大一个数量级，也就是 $B$ 中大多数元素都要排除。
	
	\item 我们假定，生成集合 $B$ 的元素比直接生成集合 $A$ 的元素容易。如果即便如此，生成 $B$ 的元素仍有困难，就可以重复这种思路，再用一次技巧，寻找一个包含 $B$ 作为子集的更大配置集合 $C$，依此类推。（细心的读者会看到，本例中我们就将这样做。）
\end{enumerate}

以上勾勒的是一个非常通用的方法，适用于许多截然不同的问题。面对具体问题，也就是具体的集合 $A$，当然需要决定：选择什么作为集合 $B$？

乐观时，也许会觉得这很容易，因为可以考虑一种技术：列出集合 $A$ 的元素必须满足的所有相互独立条件，再省略其中一个。有时它管用，但作为一般技术，却太天真了。盲目把它用于八皇后问题，其不足就会显现。我们可以用两个条件刻画解：

\begin{enumerate}[wide, nosep, label=(\arabic*)]
	\item 棋盘上有 8 个皇后。
	\item 任意两个皇后都不能互相攻击。
\end{enumerate}

省略其中任一条件，就给出集合 $B$ 的两个方案：

\begin{enumerate}[leftmargin=*, itemindent=1.9em, label=B\arabic*:]
	\item \label{en:b1-eight-queens}
	棋盘上有 $N$ 个皇后，且任意两个都不能互相攻击的全部配置。
	\item \label{en:b2-eight-queens}
	棋盘上有 8 个皇后的全部配置。
\end{enumerate}

但这两个集合都大得荒唐，会得到完全不切实际的算法。因此，必须更聪明些。迫切的问题是：“怎样更聪明？”

到这个阶段，我们有点不知所措，关心的与其说是最终程序的效率，不如说是自己思维过程的效率！因此，如果决定列举解的性质，希望找到有用线索，这是一项方向很不明确的搜索，不应投入过多脑力。也就是说，开始时应把自己限定在\textit{显然}的性质上。

（我把这道谜题作为让人清醒一下的练习，交给大学数学系的一位教员，因为他说程序设计很容易。他违背了上述规则；除了是纯数学家之外，也许还是个不高明的数学家：他开始寻找有趣而不显然的性质。他猜想，若把棋盘分为四个 $4\times 4$ 区域，每个区域必有两个皇后；然后，在尚未确信这个猜想有用之前，就开始证明它。他至今仍未解决问题；据我所知，也还没发现猜想是错的。）

好了，继续前进，把想到的显然性质列出来。

\begin{enumerate}[wide, nosep, label=(\alph*)]
	\item \label{en:property-a}
	每行至多一个皇后，需要放置 8 个皇后，而棋盘恰有 8 行。因此，每行恰有一个皇后。
	\item 同理，每列恰有一个皇后。
	\item 有 15 条“上升”对角线，每条至多一个皇后；即 8 条上升对角线恰有一个皇后，另外 7 条上升对角线为空。
	\item 同理，8 条下降对角线恰有一个皇后，另外 7 条为空。
	\item 给定任意非空皇后配置，且其中任意两个都不能互相攻击，那么移除任意一个皇后，所得配置仍具有这一性质。
\end{enumerate}

最后一条性质非常重要。（老实说，到这里，我感到无法让发明的突然一跃变得平缓！）用前面的术语说，它告诉我们，集合\hyperref[en:b1-eight-queens]{B1}中的任意非空配置都具有这个特点。如果从一个解——即集合\hyperref[en:b1-eight-queens]{B1}中的八皇后配置——出发，移除一个皇后，就得到集合\hyperref[en:b1-eight-queens]{B1}中的七皇后配置；再移除一个，仍得到\hyperref[en:b1-eight-queens]{B1}中的配置；可以一直重复，直到棋盘为空。我们可以把这个过程拍成电影；倒放时，就会看见：从空棋盘开始，经过集合\hyperref[en:b1-eight-queens]{B1}中的配置，每次加入一个皇后，逐步建立那个解。（倒放电影的技巧对读者是否有帮助，不由我判断；我提到它，只因为知道这类办法对我有帮助。）拍电影时，任意一个解都可以通过许多方式化为空棋盘；而倒放时，恰好也有同样多的方式把这个解建立起来。能利用这种自由吗？我们已经因为集合\hyperref[en:b1-eight-queens]{B1}太大而舍弃它，但也许可以找到一个合适子集，使其中每个非空配置，都是子集中唯一一个其他配置添加一个皇后所得的扩展。“扩展性质”表明，我们愿意考虑棋盘上少于 8 个皇后的配置，并希望通过向现有配置加入一个皇后，形成新配置；这想必是一项相对简单的操作。于是，注意力立即转到这个仍然神秘的集合 $B$ 的元素\textit{生成}上。例如，按什么顺序生成？这又引出一个至今完全没有考虑的问题：按什么顺序生成解，也就是集合 $A$ 的元素？能否提出合理建议，希望从中找到线索？（根据我的经验，这类关于顺序的问题，通常很有启发。原因不只在于：我们要写顺序程序，按定义它会以某个顺序生成解，因此总要决定这个顺序。顺序的决定通常还能提供线索，证明程序会生成\textit{全部}解，并且每个解只生成\textit{一次}。）

在此之前，应当问：有了解之后，怎样刻画它？必须给出 8 个皇后的位置。皇后本身没有顺序，但行列有；可以假定它们都从 0 到 7 编号。性质 \ref{en:property-a} 告诉我们，每行恰有一个皇后，因此可以按所在行的编号为皇后排序。于是，每个八皇后配置，都可由\textbf{integer array} $x[0:7]$ 的值给出，其中

\noindent
$x[i] =$ 第 $i$ 行皇后所在列的编号。

每个解便是一个“八位数词”——$x[0] \dots x[7]$。我能想到唯一合理的生成顺序，就是字典序。

\noindent
\textit{注}。这样，我们就为行列处理方式不同的算法开了门，而原问题对行列是对称的！考虑不对称算法，恰恰是上述思考教给我们的东西！

再回到字典序：现在已接近熟悉的地方。如果集合 $A$ 的元素要按字典序生成，而且要从较大集合 $B$ 中筛选得到，那么标准技术就是：同样按字典序生成 $B$ 的元素，再按子集元素在 $B$ 中出现的次序输出它们。

首先，生成所有 $x[0] = 0$ 的解——如果有；然后生成 $x[0] = 1$ 的解——如果有；依此类推。在 $x[0]$ 固定的那些解中，应先生成 $x[1] = 0$ 的解，再生成 $x[1] = 1$ 的解，依此类推。换句话说，第 0 行皇后先放在第 0 列——比如左下角的格子——并保持在那里，直到集合 $A$（和 $B$）中第 0 个皇后处于该位置的所有元素都生成完毕，再向右移动一格到下一列。对于第 0 个皇后的每个位置，第 1 个皇后从左到右走过第 1 行，跳过被第 0 个皇后控制的格子；对于前两个皇后的每种组合位置，第 2 个皇后从左到右走过第 2 行，跳过所有被前面皇后控制的格子；依此类推。

现在，集合 $B$ 找到了！它确实是\hyperref[en:b1-eight-queens]{B1}的子集：$B$ 包含所有这样的配置——前 $N$ 行每行一个皇后，且任意两个皇后都不能互相攻击。

判断 $B$ 的元素是否也属于 $A$，判据就是 $N = 8$。

确定 $B$ 的选择之后，任务变成按字典序生成它的元素。可以尝试用操作“生成 $B$ 的下一个元素”，写成以下程序：

初始化空棋盘；

\textbf{\textbf{repeat}} 生成 $B$ 的下一个元素；

\quad \textbf{if} $N = 8$ \textbf{then} 打印配置

\textbf{until} $B$ 已穷尽。

\noindent
（这里利用了以下事实：空棋盘属于 $B$，不属于 $A$，而且不是 $B$ 的唯一元素。并未对解的存在作任何假定。）

但这种结构有两个理由不太理想。第一，没有现成判据，能在遇到 $B$ 最后一个元素时认出它；很可能必须推广“生成 $B$ 的下一个元素”，让它用于 $B$ 的最后一个“真正”元素时，产生“$B$ 已穷尽”的指示。第二，怎样构造“生成 $B$ 的下一个元素”并不明显：棋盘上皇后的个数可能不变、减少，也可能增加。

所以，这不太理想。能怎么办？只要把集合 $B$ 的配置序列看作一个单一、没有变化的序列，不将它分为先后相继的子序列，相应程序结构就会像刚才那样，是一个循环。如果要寻找别的程序结构，就\textit{必须}问自己：“怎样把集合 $B$ 的配置序列，分组为先后相继的子序列？”

想到 $B$ 的配置必须按字典序生成，再想到词典的主要分组方式——按首字母——第一种分组便很明显：按第 0 个皇后的位置。

于是，生成集合 $B$ 的全部元素——暂且不管打印其中也属于 $A$ 的配置——表现为

初始化空棋盘；

$h \coloneq 0$;

\textbf{repeat} 在格子 $[0,\ h]$ 上放置皇后；

\quad 生成第 0 个皇后位置

\quad\quad 固定的全部配置；

\quad 从格子 $[0,\ h]$ 移除皇后；

\quad $h \coloneq h + 1$

\textbf{until} $h = 8$。

但问题再次出现：怎样对第 0 个皇后固定的所有配置分组？已经给出答案：按第 1 个皇后所在列的编号递增分组，也就是

$h1 \coloneq 0$;

\textbf{repeat if} 格子 $[1,\ h1]$ 可用 \textbf{do}

\quad \textbf{begin} 在格子 $[1,\ h1]$ 上放置皇后；

\quad \quad 生成前两个皇后

\quad \quad \quad 位置固定的全部配置；

\quad \quad 从格子 $[1,\ h1]$ 移除皇后

\quad \textbf{end};

\quad $h1 \coloneq h1 + 1$

\textbf{until} $h1 = 8$。

对于“生成前两个皇后位置固定的全部配置”，也可以写出类似程序，依此类推。把它们互相嵌入，就得到一个有八层嵌套循环的正确程序；但这些循环都非常、非常相似。这样做有两个缺点：

\begin{enumerate}[wide, nosep, label=(\arabic*)]
	\item 需要写大量繁琐文字。
	\item 得到的程序只能解决 $8\times 8$ 棋盘的问题；如果要解 $10\times 10$ 棋盘的同一道谜题，就需要一个新的、更长的程序。
\end{enumerate}

我们想找到一种方法，使所有循环都由同一段程序文本控制执行。能否让循环文本相同？能否利用这种相同性？

首先，最外层和最内层循环都是例外。

最外层循环的例外在于：它不检查 square$[0,\ h]$ 是否可用，因为我们知道它可用。但正因为知道它可用，加入条件子句并无妨碍：

\quad \textbf{if} 格子 $[0,\ h]$ 可用 \textbf{do}

\noindent
这样，最外层循环就与接下来六层具有相同模式。

最内层循环的例外在于：棋盘上放满 8 个皇后后，再生成这些皇后固定的全部配置已没有意义，因为棋盘满了。此时应打印配置，因为已经找到一个同时属于集合 $B$ 与集合 $A$ 的元素。把“生成”那一行替换为以下内容，就能让最内层循环与外面的七层对应起来：

\textbf{if} 棋盘已满 \textbf{\textbf{then}} 打印配置

\quad \textbf{else} 生成扩展当前配置的

\quad\quad 全部配置

为此，引入一个全局变量，比如“$n$”，计数当前棋盘上的皇后。“棋盘已满”的检查变为“$n = 8$”；格子操作也可以用“$n$”作第一下标。

现在，八个循环的唯一差别，是每个都有“自己的 $h$”。到这一阶段，对于能否利用循环的相同性，已经可以肯定回答。通过一个递归过程，比如“generate”，只把循环描述一次，就能引发八个嵌套循环的执行顺序。用它后，程序本身简化为

\quad 初始化空棋盘；$n \coloneq 0$;

\quad generate

\noindent
而“generate”递归定义如下：

\textbf{procedure} generate;

\textbf{begin integer} $h$;

$h \coloneq 0$;

\textbf{repeat if} 格子 $[n,\ h]$ 可用 \textbf{do}

\quad \textbf{\textbf{begin}} 在格子 $[n,\ h]$ 上放置皇后；$n \coloneq n + 1$;

\quad \quad \textbf{if} $n = 8$ \textbf{then} 打印配置

\quad \quad \textbf{else} generate;

\quad \quad $n \coloneq n - 1$; 从格子 $[n,\ h]$ 移除皇后

\quad \textbf{end};

\quad $h \coloneq h + 1$

\quad \textbf{until} $h = 8$

\textbf{end}

“generate”的每次激活，都会引入它私有的局部变量 $h$，从而提供写八层嵌套循环时所需的 $h,\ h1,\ \dots,\ h8$。

虽然程序在这一细节层次上正确，但还未完成，即尚未细化到程序设计语言所要求的标准细节程度。下一步，应决定用什么约定表示棋盘配置。我们已经多少决定，采用

\quad \textbf{integer array} $x[0:7]$

\noindent
按次序给出各皇后所在列的编号，并且采用

\quad \textbf{integer} $n$

\noindent
表示棋盘上的皇后个数。更精确地说，

$n =$ 棋盘上的皇后个数

\hangindent=2\parindent 对 $0 \leqslant i < n$，$x[i] =$ 第 $i$ 行皇后所在列的编号。

数组 $x$ 和标量 n，合起来已足以确定集合 $B$ 中的任意配置，而棋盘上只会出现这些配置。因此，在\textit{逻辑上}，无需更多变量；但我们仍将引入几个，因为从实践上看，它们可以很好地发挥作用。问题是：只用上述信息，频繁检查下一空行中的某个给定格子是否不受攻击，会很麻烦、费时。这里要用“\nameref{sec:trading-storage-for-computation-speed}”一节描述的标准技术（参见第 \pageref{sec:trading-storage-for-computation-speed} 页）。已存储的自变量，在这里是棋盘上的皇后配置；这个值并不剧烈变化——当然不会——我们只会加入或移除一个皇后。因此，寻找一些额外的表：内容是当前配置的函数，可以帮助判断格子是否可用；而且加入或移除一个皇后时，容易更新。

怎样做？可以考虑一个 $8\times 8$ 的布尔数组，表示每个格子是否可用。如果为整个棋盘这样做，加入一个皇后可能涉及处理 28 个格子。但移除皇后就很麻烦，因为不再受\textit{她}攻击的格子，未必真的可用：还可能受剩下的一个或几个皇后攻击。对此，有一个同样标准的补救：每格关联的不是布尔变量，而是整数计数器，计数有多少皇后控制该格。加入皇后时，给至多 28 个计数器加 1；移除时减 1；相应计数器为零，格子就可用。我们可以这样做，但要问是否过头：放置或移除一个皇后，附带 28 次调整，开销确实很大。

每个我们关心其可用性的格子，都涉及一行（按定义这一行为空，所以不必管）、8 列之一（必须仍为空）、15 条上升对角线之一（必须仍为空），以及 15 条下降对角线之一（必须仍为空）。这提示我们记录

\begin{enumerate}[wide, nosep, label=(\arabic*)]
	\item 哪些列可用；
	\item 哪些上升对角线可用；
	\item 哪些下降对角线可用。
\end{enumerate}

每列或每条对角线只被一个皇后占据，因此不用计数器，布尔变量就够了。列可以直接用列编号标识，于是引入

\textbf{boolean array} col$[0:7]$

\noindent
其中，“col$[i]$”表示第 $i$ 列仍可用。

怎样标识对角线？沿上升对角线，行编号减列编号的差不变；沿下降对角线，两者的和不变。因此，差与和分别是最容易用来区分这两类对角线的下标，于是引入

\quad \textbf{boolean array} up$[-7:+7]$, down$[0:14]$

\noindent
记录哪些对角线可用。

判断 square$[n,\ h]$ 是否可用，变为

\quad col$[h]$ \textbf{and} up$[n-h]$ \textbf{and} down$[n+h]$，

\noindent
放置和移除皇后，都只需调整三个布尔值，每个数组一个。

最终程序引入变量“$k$”，作一般计数用；语句和表达式都加了标号（大写字母）。请注意，我们合并了两个描述层次：上层中的语句和函数，现在变成说明性标号。

给出最终程序后，最后一节也就结束了。我们试图展示：怎样通过一种推理方式，发现回溯技术；又怎样通过一种推理方式，发现描述它的递归过程。也许本节最重要的启示是：全部分析与综合，都能在决定配置如何在机器内表示——以及采用多少冗余信息——之前完成。确实，只有最终找到方便的配置表示，这种思考才会结出成果。但在思维上划出一个抽象层次，允许自己暂时不操心表示问题，似乎至关重要。

最后，我感谢一直读到这里的读者，感谢你的耐心。
\bigskip

\noindent
\textbf{begin} \textbf{integer} $n$, $k$; \textbf{integer array} $x[0:7]$;

\noindent
\quad \quad \textbf{boolean array} col$[0:7]$, up$[-7:+7]$, down$[0:14]$;

\noindent
\quad \textbf{procedure} generate;

\noindent
\quad \textbf{begin integer} $h$;

\noindent
\quad \quad $h \coloneq 0$;

\noindent
\quad \quad \textbf{repeat}

\noindent
\quad \quad \quad \textbf{if} 格子 $[n,h]$ 可用：(col$[h]$ \textbf{and} up$[n-h]$ \textbf{and} down$[n+h]$) \textbf{do}

\noindent
\quad \quad \quad \textbf{begin} 在格子 $[n,h]$ 放置皇后：

\noindent
\quad \quad \quad \quad \quad $x[n] \coloneq h$; col$[h]\coloneq$ \textbf{false}; up$[n-h]\coloneq$ \textbf{false};

\noindent
\quad \quad \quad \quad \quad down$[n+h]\coloneq$ \textbf{false}; $n\coloneq n + 1$;

\noindent
\quad \quad \quad \quad \textbf{if} 棋盘已满：$(n = 8)$ \textbf{then}

\noindent
\quad \quad \quad \quad \textbf{begin} 打印配置：

\noindent
\quad \quad \quad \quad \quad $k\coloneq 0$; \textbf{repeat} print$(x[k])$; $k\coloneq k + 1$ \textbf{until} $k = 8$; newline


\noindent
\quad \quad \quad \quad \textbf{end}

\noindent
\quad \quad \quad \textbf{else} generate;

\noindent
\quad \quad \quad \quad $n \coloneq n - 1;$

\noindent
\quad \quad \quad \quad 从格子 $[n,h]$ 移除皇后：

\noindent
\quad \quad \quad \quad \quad down$[n+h]\coloneq$ \textbf{true}; up$[n-h] \coloneq$ \textbf{true}; col$[h]\coloneq$ \textbf{true}

\noindent
\quad \quad \quad \textbf{end};

\noindent
\quad \quad \quad $h\coloneq h + 1$

\noindent
\quad \quad \textbf{until} $h = 8$

\noindent
\quad \textbf{end};

\noindent
\quad 初始化空棋盘：

\noindent
\quad \quad $n \coloneq 0$;

\noindent
\quad \quad $k\coloneq 0$; \textbf{repeat} col$[k] \coloneq$ \textbf{true}; $k\coloneq k + 1$ \textbf{until} $k = 8$;

\noindent
\quad \quad $k\coloneq 0$; \textbf{repeat} up$[k-7] \coloneq$ \textbf{true}; down$[k] \coloneq$ \textbf{\textbf{true}}; $k\coloneq k + 1$ \textbf{until} $k=15$;

\noindent
\quad generate

\noindent
\textbf{end}
